\documentclass[12pt, a4paper]{report}

% --- Pacotes Essenciais ---
\usepackage[utf8]{inputenc}
\usepackage[portuguese]{babel}
\usepackage{amsmath, amssymb, amsfonts, amsthm}
\usepackage{mathtools}
\usepackage{graphicx}
\usepackage{geometry}
\usepackage{hyperref}
\usepackage{booktabs}
\usepackage{tabularx}
\usepackage{enumitem}
\usepackage{listings}
\usepackage{multicol}   % para listas a duas colunas

% --- Pacotes para Estilização Avançada ---
\usepackage[dvipsnames]{xcolor}
\usepackage[most]{tcolorbox}
\tcbuselibrary{breakable, skins}
\usepackage{titlesec}
\usepackage{tikz}
\usetikzlibrary{arrows.meta, positioning, calc}
\usepackage{pgfplots}
\pgfplotsset{compat=1.18}

% Configuração das Margens
\geometry{left=2.5cm, right=2.5cm, top=3cm, bottom=3cm}

% Definição de Cores Profissionais
\definecolor{primary}{RGB}{0,72,144}
\definecolor{accent}{RGB}{220,50,47}
\definecolor{soft}{RGB}{240,246,255}
\definecolor{softgray}{RGB}{245,245,245}
\definecolor{warnyellow}{RGB}{255,248,220}
\definecolor{greenbg}{RGB}{235,255,235}
\definecolor{redbg}{RGB}{255,235,235}

% Configuração do Hyperref
\hypersetup{
    colorlinks=true,
    linkcolor=primary,
    urlcolor=primary,
    citecolor=primary
}

% Formatação Elegante dos Capítulos e Secções
\titleformat{\chapter}[display]
  {\normalfont\Huge\bfseries\color{primary}}{\chaptertitlename\ \thechapter}{10pt}{\Huge}[\vspace{1ex}\color{primary}\titlerule]
\titleformat{\section}
  {\Large\bfseries\color{primary}}{\thesection.}{0.5em}{}
\titleformat{\subsection}
  {\large\bfseries\color{primary!80!black}}{\thesubsection.}{0.5em}{}

% Estilos Customizados de Caixas (tcolorbox)
\tcbset{
  basestyle/.style={
    breakable, enhanced, fonttitle=\bfseries, coltitle=white,
    attach boxed title to top left={yshift=-2mm, xshift=4mm},
    boxed title style={rounded corners, size=small},
  }
}
\newtcolorbox{defbox}[1][]{basestyle, colback=soft, colframe=primary, title={Definição}, boxed title style={colback=primary}, #1}
\newtcolorbox{thmbox}[1][]{basestyle, colback=greenbg, colframe=green!60!black, title={Teorema}, boxed title style={colback=green!60!black}, #1}
\newtcolorbox{exbox}[2][]{basestyle, colback=softgray, colframe=gray!60, title={Exemplo~#2}, boxed title style={colback=gray!60}, #1}
\newtcolorbox{solbox}[1][]{breakable, enhanced, colback=white, colframe=primary!40, title={\textbf{Resolução}}, fonttitle=\bfseries\color{primary}, borderline west={3pt}{0pt}{primary}, #1}
\newtcolorbox{warnbox}[1][]{breakable, enhanced, colback=warnyellow, colframe=orange!80!black, title={\textbf{Erro Comum / Aviso}}, fonttitle=\bfseries\color{orange!80!black}, borderline west={4pt}{0pt}{orange!80!black}, #1}
\newtcolorbox{notebox}[1][]{breakable, enhanced, colback=softgray, colframe=gray!60, title={\textbf{Nota}}, fonttitle=\bfseries, borderline west={4pt}{0pt}{gray!60}, #1}
\newtcolorbox{exercisebox}[2][]{basestyle, colback=redbg!40!white, colframe=accent, title={Exercício Prático~#2}, boxed title style={colback=accent}, #1}
\newtcolorbox{oralbox}[1][]{basestyle, colback=primary!8, colframe=primary!70, title={Questão de Defesa Oral}, boxed title style={colback=primary!70}, #1}
\newtcolorbox{tipbox}[1][]{basestyle, colback=yellow!10!white, colframe=orange!80!black, title={Dica / Sugestão}, boxed title style={colback=orange!80!black}, #1}

\lstset{
    language=Matlab,
    basicstyle=\ttfamily\footnotesize,
    keywordstyle=\color{blue}\bfseries,
    commentstyle=\color{teal},
    stringstyle=\color{red},
    numbers=left,
    numberstyle=\tiny\color{gray},
    stepnumber=1,
    breaklines=true,
    frame=single,
    backgroundcolor=\color{gray!5},
    captionpos=b,
    showstringspaces=false
}

% Macros úteis
\newcommand{\dd}{\mathrm{d}}

% --- Informações da Capa ---
\title{
    \vspace{-2cm}
    \textbf{Análise Matemática}\\
    \large Relatório Teórico
}
\author{
    David Domingues (2022125309) \\
    Diana Gomes (2019101181) \\
    Guilherme Baltazar (2025209893) \\
    Joana Venâncio (2016123347) \\
    Leonardo Costa (2023137263) \\
    Marcelo Marques (2022102560) \\
    Miguel Reis (2024164633) \\
}
\date{Politécnico de Leiria \\ 2026}

% =========================================================
% O FILTRO: Aqui dizes ao Overleaf o que queres compilar HOJE
% =========================================================
\includeonly{
    %AM - Report/Chapter 0/0_1 Teoria dos Conjuntos,
    %AM - Report/Chapter 0/0_2 Equivalence Relations,
    %AM - Report/Chapter 0/0_3 Espacos Metricos
    %AM - Report/Chapter 0/0_4 Espacos Topologicos,
    %AM - Report/Chapter 0/0_5_Limites_Continuidade,
    %AM - Report/Chapter 0/0_6 Perspetiva Estrutural
    %AM - Report/Chapter 1/1_1_Conceitos_de_Derivada,
    %AM - Report/Chapter 1/1_2_Derivadas_Especiais
    %AM - Report/Chapter 1/1_3_Derivadas_de_Ordem Superior,
    %AM - Report/Chapter 1/1_4_Curvas_Parametricas
    %AM - Report/Chapter 2/2_1_Integrais_Indefinidas
    %AM - Report/Chapter 2/2_2_Integrais_Definidas,
    %AM - Report/Chapter 2/2_3_Teorema_Fundamental
    %AM - Report/Chapter 2/2_4_Aplicacoes_E_2_5_Comprimento_Arco
    %AM - Report/Chapter 3/3_1_Definicoes_Basicas
    %AM - Report/Chapter 3/3_2_Limites_e_Continuidade
    %AM - Report/Chapter 3/3_3_Derivadas_Parciais_Diferenciabilidade
    %AM - Report/Chapter 3/3_4_ Extremos_Livres
    %AM - Report/Chapter 4/4_ Integrais_Duplas
    AM - Report/Chapter 5/5_ODEs
}

\begin{document}

\maketitle
\tableofcontents
\newpage

% ---------------------------------------------------------
% LISTA COMPLETA: Não comentar estas linhas! 
% O \includeonly lá em cima é que decide quais aparecem no PDF.
% ---------------------------------------------------------

% ---------------------------------------------------------
% CAPÍTULO 0
% ---------------------------------------------------------
% ════════════════════════════════════════════════════════════════
% 0.1 - TEORIA DE CONJUNTOS
% ════════════════════════════════════════════════════════════════

\chapter{Conceitos Matemáticos Fundamentais}

% ────────────────────────────────────────────────────────
\section{Teoria de Conjuntos}

% ────────────────────────────────────────────────────────
\subsection{Enquadramento Teórico}

A teoria de conjuntos de Neumann--Bernays--Gödel (NBG) é uma formalização axiomática que resolve paradoxos clássicos como o de Russell, distinguindo \emph{classes} de \emph{conjuntos}. 

\begin{notebox}[title={Contexto Histórico e Comparação NBG vs.\ ZFC}]
  \begin{itemize}
    \item \textbf{ZFC} (Zermelo--Fraenkel + Axioma da Escolha): trabalha só com conjuntos e restringe a formação de conjuntos através de axiomas.
    \item \textbf{NBG} (von Neumann--Bernays--Gödel): introduz \emph{conjuntos} e \emph{classes}; é conservativo sobre ZFC, mas permite falar de classes próprias (como a classe de todos os conjuntos) sem inconsistência.
  \end{itemize}
\end{notebox}

\begin{defbox}[title={Classes e Conjuntos}]
  \begin{itemize}
      \item \textbf{Classe:} Coleção de objetos que partilham uma propriedade lógica específica $C = \{ x \mid \varphi(x) \}$.
      \item \textbf{Conjunto:} É uma classe $X$ tal que existe uma classe $Y$ com $X \in Y$.
      \item \textbf{Classe própria:} É uma classe que \emph{não} é um conjunto (não pertence a nenhuma outra classe). Ex: Classe universal $V$, Classe de todos os ordinais $\mathrm{Ord}$.
  \end{itemize}
\end{defbox}

\begin{thmbox}[title={O Paradoxo de Russell}]
  A coleção $R = \{x \mid x \notin x\}$ não pode ser um conjunto. Se $R$ fosse um conjunto, teríamos a contradição: $R \in R \iff R \notin R$. Portanto, $R$ é uma classe própria.
\end{thmbox}

\begin{defbox}[title={Subconjuntos e Operações}]
  \begin{itemize}
    \item \textbf{Conjunto das partes:} $\mathcal{P}(A) = \{X \mid X \subseteq A\}$.
    \item \textbf{União:} $A \cup B = \{x \mid x \in A \lor x \in B\}$
    \item \textbf{Interseção:} $A \cap B = \{x \mid x \in A \land x \in B\}$
    \item \textbf{Diferença:} $A \setminus B = \{x \mid x \in A \land x \notin B\}$
    \item \textbf{Leis de De Morgan:} $(A \cup B)^c = A^c \cap B^c$ e $(A \cap B)^c = A^c \cup B^c$.
  \end{itemize}
\end{defbox}

% ────────────────────────────────────────────────────────
\subsection{Exemplos Resolvidos}

\begin{exbox}{1 --- Identificação e Natureza das Classes}
  Determine se a classe $B = \{x \mid x \text{ é um número primo}\}$ é um conjunto ou classe própria.
\end{exbox}
\begin{solbox}
  É um \textbf{conjunto}. Como é um subconjunto de $\mathbb{N}$ (garantido pelo Axioma da Separação) e $\mathbb{N}$ é um conjunto pelo Axioma do Infinito, $B$ é suficientemente pequeno para ser um conjunto.
\end{solbox}

\begin{exbox}{2 --- Operações e De Morgan}
  Seja $U = \{1,2,3,4,5,6\}$, $A = \{1,2,3\}$ e $B = \{3,4,5\}$. Calcule $(A \cup B)^c$ e verifique De Morgan.
\end{exbox}
\begin{solbox}
  $A \cup B = \{1,2,3,4,5\} \implies (A \cup B)^c = \{6\}$. \\
  Por outro lado, $A^c = \{4,5,6\}$ e $B^c = \{1,2,6\}$. A interseção $A^c \cap B^c = \{6\}$. A lei verifica-se. \checkmark
\end{solbox}

% ────────────────────────────────────────────────────────
\subsection{Exercícios Práticos}

\begin{exercisebox}{1}
  Dados $U=\{1,2,\ldots,10\}$, $A=\{1,3,5,7,9\}$ e $B=\{1,2,3,5\}$, determine a diferença simétrica $A \triangle B$.
\end{exercisebox}
\begin{solbox}
  $A \setminus B = \{7,9\}$ e $B \setminus A = \{2\}$. \\
  A diferença simétrica é $(A \setminus B) \cup (B \setminus A) = \{2,7,9\}$.
\end{solbox}

% ────────────────────────────────────────────────────────
\subsection{Questões Propostas para Defesa Oral}

\begin{oralbox}
  \textbf{Q1.} Na teoria NBG, explique a diferença entre \emph{Conjunto} e \emph{Classe}. Dê um exemplo de uma classe própria e justifique por que razão não pode ser um conjunto.
\end{oralbox}

\begin{oralbox}
  \textbf{Q2.} Enuncie o Paradoxo de Russell. Qual é a contradição que ele revela na teoria ingénua de conjuntos e como é que a teoria NBG o resolve?
\end{oralbox}
% ════════════════════════════════════════════════════════════════
% 0.2 - RELAÇÕES DE EQUIVALÊNCIA E CLASSES DE EQUIVALÊNCIA
% ════════════════════════════════════════════════════════════════

\section{Relações de Equivalência e Classes de Equivalência}

% ────────────────────────────────────────────────────────
\subsection{Enquadramento Teórico}

Uma \textbf{relação} num conjunto $A$ dita as regras que definem quando dois elementos desse conjunto estão relacionados.

\begin{defbox}[title={Relação de Equivalência}]
  Uma \textbf{relação de equivalência} (frequentemente denotada por $\sim$) é uma relação especial que tem de satisfazer, em simultâneo, três propriedades fundamentais para quaisquer elementos $a, b, c \in A$:
  \begin{enumerate}
      \item \textbf{Reflexividade:} $a \sim a$. (Todo o elemento está relacionado consigo mesmo).
      \item \textbf{Simetria:} Se $a \sim b$, então $b \sim a$. (A relação é recíproca).
      \item \textbf{Transitividade:} Se $a \sim b$ e $b \sim c$, então $a \sim c$. (A relação transfere-se em cadeia).
  \end{enumerate}
\end{defbox}

\begin{warnbox}[title={Nota Importante}]
  As relações de equivalência só se definem entre elementos do mesmo conjunto. Basta que \textbf{uma} destas três propriedades falhe para que a relação deixe imediatamente de ser de equivalência.
\end{warnbox}

Quando temos uma relação de equivalência $\sim$ num conjunto $A$, esta divide o conjunto em subgrupos mais pequenos.

\begin{defbox}[title={Classes de Equivalência e Conjunto Quociente}]
  A \textbf{classe de equivalência} de um elemento $a$, denotada por $[a]$ ou $\overline{a}$, é o grupo de elementos de $A$ que se relacionam com $a$ segundo a regra definida:
  \[
      [a] = \{ x \in A \mid x \sim a \}
  \]
  A união de todas as classes de equivalência constitui o conjunto inteiro, e classes distintas não têm elementos em comum (são disjuntas). Isto significa que a relação cria uma \textbf{partição} do conjunto $A$. 
  
  Ao conjunto formado por todas as classes de equivalência distintas chamamos \textbf{conjunto quociente}, denotado por $A/\sim$.
\end{defbox}

% ────────────────────────────────────────────────────────
\subsection{Exemplos Resolvidos}

\begin{exbox}{1 --- Relação Numérica em Conjunto Finito}
  Considere o conjunto $A = \{2, 4, 6, 8\}$ e a relação $a \sim b$ definida por: \textit{``$a$ e $b$ estão relacionados se a diferença $a - b$ for um múltiplo de 2''}. Verifique se é uma relação de equivalência.
\end{exbox}
\begin{solbox}
  Vamos verificar as propriedades para elementos arbitrários:
  \begin{itemize}
      \item \textbf{Reflexiva:} $2 \sim 2$, porque $2 - 2 = 0$, e $0$ é múltiplo de $2$. \checkmark
      \item \textbf{Simétrica:} Se $2 \sim 4$, então $4 - 2 = 2$ (múltiplo de 2). Logo, $4 \sim 2$. \checkmark
      \item \textbf{Transitiva:} Se $2 \sim 4$ (pois $2-4=-2$) e $4 \sim 6$ (pois $4-6=-2$), então $2 \sim 6$, porque $2 - 6 = -4$, que é também múltiplo de $2$. \checkmark
  \end{itemize}
  Como as três propriedades se verificam, trata-se de facto de uma relação de equivalência.
\end{solbox}

\begin{exbox}{2 --- O Percurso na Disciplina de PAM}
  Considere o conjunto de alunos: \\
  $Grupo = \{\text{David, Diana, Guilherme, Joana, Leonardo, Marcelo, Miguel}\}$. \\
  Definimos a relação de equivalência $\sim$: \textit{``Dois alunos relacionam-se se têm o mesmo histórico face à cadeira de PAM (ambos a fizeram ou ambos não a fizeram)''}. Determine as classes de equivalência e o conjunto quociente.
\end{exbox}
\begin{solbox}
  Sabendo que David, Guilherme, Joana e Marcelo vieram de PAM, enquanto Diana, Leonardo e Miguel não vieram, a relação particiona o grupo.

  \textbf{Classes de Equivalência:} \\
  Podemos representar cada classe usando um dos seus elementos (o ``representante'' da classe):
  \begin{itemize}
      \item $[\text{David}] = \{\text{David, Guilherme, Joana, Marcelo}\}$ (Classe dos que fizeram PAM)
      \item $[\text{Diana}] = \{\text{Diana, Leonardo, Miguel}\}$ (Classe dos que não fizeram PAM)
  \end{itemize}
  Note-se que podíamos escrever $[\text{Guilherme}]$ em vez de $[\text{David}]$, pois $[\text{Guilherme}] = [\text{David}]$. 
  
  O \textbf{conjunto quociente} é o conjunto destas duas classes:
  \[
      Grupo/\sim \; = \{ [\text{David}], [\text{Diana}] \}
  \]
\end{solbox}

\begin{exbox}{3 --- Resto da Divisão por 3 no Conjunto $\mathbb{Z}$}
  Considere o conjunto infinito dos números inteiros ($\mathbb{Z}$) e a relação $a \sim b$ definida por: \textit{``$a$ e $b$ têm o mesmo resto quando divididos por 3''}. Matematicamente, isto equivale a dizer que a diferença $a - b$ é um múltiplo de 3 (relação conhecida como congruência módulo 3, $a \equiv b \pmod{3}$). Determine as classes de equivalência.
\end{exbox}
\begin{solbox}
  Ao dividirmos qualquer número inteiro por 3, os únicos restos possíveis são 0, 1 ou 2. Por isso, esta relação gera exatamente três classes de equivalência distintas, particionando o infinito conjunto $\mathbb{Z}$:
  \begin{itemize}
      \item $[0] = \{\dots, -6, -3, 0, 3, 6, 9, \dots\}$ (Números com resto 0; múltiplos de 3)
      \item $[1] = \{\dots, -5, -2, 1, 4, 7, 10, \dots\}$ (Números com resto 1)
      \item $[2] = \{\dots, -4, -1, 2, 5, 8, 11, \dots\}$ (Números com resto 2)
  \end{itemize}
  Qualquer outro inteiro recairá numa destas classes. Por exemplo, o número $4$ pertence à classe do $1$, logo $[4] = [1]$. 
  
  O \textbf{conjunto quociente} para esta relação é formado por estas três classes:
  \[
      \mathbb{Z}/\sim \; = \{ [0], [1], [2] \}
  \]
\end{solbox}

% ────────────────────────────────────────────────────────
\subsection{Exercícios Práticos}

\begin{exercisebox}{1}
  Considere o conjunto de alunos $A = \{\text{Ana, João, Marta, Pedro}\}$. Sabe-se que a Ana e o João nasceram no ano 2000, e a Marta e o Pedro nasceram em 2001. Defina a relação $x \sim y \iff \text{``x e y têm o mesmo ano de nascimento''}$.
  \begin{enumerate}
      \item Mostre, usando a definição formal, que esta é uma relação de equivalência.
      \item Determine as classes de equivalência e o conjunto quociente $A/\sim$.
  \end{enumerate}
\end{exercisebox}
\begin{solbox}
  \textbf{1. Verificação das Propriedades:}
  \begin{itemize}
      \item \textbf{Reflexividade:} Qualquer pessoa $x$ nasceu no mesmo ano que si própria ($x \sim x$). \checkmark
      \item \textbf{Simetria:} Se a pessoa $x$ nasceu no mesmo ano que $y$, então $y$ nasceu no mesmo ano que $x$ ($x \sim y \implies y \sim x$). \checkmark
      \item \textbf{Transitividade:} Se $x$ nasceu no mesmo ano que $y$, e $y$ nasceu no mesmo ano que $z$, então $x$ nasceu no mesmo ano que $z$ ($x \sim y \land y \sim z \implies x \sim z$). \checkmark
  \end{itemize}
  
  \textbf{2. Classes de Equivalência e Conjunto Quociente:} \\
  As classes de equivalência agrupam os alunos pelo ano de nascimento:
  \begin{itemize}
      \item $[\text{Ana}] = \{\text{Ana, João}\}$
      \item $[\text{Marta}] = \{\text{Marta, Pedro}\}$
  \end{itemize}
  O conjunto quociente é formado por estas duas classes: 
  \[
      A/\sim \; = \{ [\text{Ana}], [\text{Marta}] \}
  \]
\end{solbox}

% ────────────────────────────────────────────────────────
\subsection{Questões Propostas para Defesa Oral}

\begin{oralbox}
  \textbf{Q1.} Usando o Exemplo 2 do relatório (o percurso dos alunos em PAM), explique por que motivo a notação $[\text{David}] = [\text{Guilherme}]$ está correta e relacione isso com a propriedade de \emph{partição} de um conjunto.
\end{oralbox}

\begin{oralbox}
  \textbf{Q2.} No Exemplo 3, se definíssemos a relação como "resto da divisão por 5" em vez de 3, quantos elementos teria o conjunto quociente $\mathbb{Z}/\sim$ e quais seriam os seus representantes mais simples?
\end{oralbox}
% ════════════════════════════════════════════════════════════════
% 0.3 - ESPAÇOS MÉTRICOS
% ════════════════════════════════════════════════════════════════

\section{Espaços Métricos}

% ────────────────────────────────────────────────────────
\subsection{Enquadramento Teórico}

A noção de distância é um dos conceitos mais intuitivos da matemática, mas para ser aplicada a conjuntos abstratos (como conjuntos de funções ou matrizes), precisa de ser formalizada. A estrutura que generaliza o conceito de "distância" chama-se \textbf{espaço métrico}.

\begin{defbox}[title={Métrica e Espaço Métrico}]
  Seja $X$ um conjunto. Uma \textbf{métrica} em $X$ é uma função
  \[
    d: X \times X \to \mathbb{R}
  \]
  tal que, para quaisquer $x,y,z \in X$, satisfaz as seguintes três propriedades:
  \begin{enumerate}[label=(\roman*)]
      \item \textbf{Positividade:} $d(x,y) \ge 0$ e $d(x,y)=0 \iff x=y$;
      \item \textbf{Simetria:} $d(x,y) = d(y,x)$;
      \item \textbf{Desigualdade Triangular:} $d(x,z) \le d(x,y) + d(y,z)$.
  \end{enumerate}
  O par $(X,d)$ chama-se \textbf{espaço métrico}.
\end{defbox}

\begin{defbox}[title={Bolas Abertas e Convergência}]
  Num espaço métrico $(X,d)$, para um ponto $x\in X$ e um raio $r>0$, define-se a \textbf{bola aberta} centrada em $x$ com raio $r$ como:
  \[
    B_r(x) = \{y\in X : d(x,y)<r\}.
  \]

  Diz-se que uma sequência $(x_n)$ \textbf{converge} para $x$ (denotado por $x_n \to x$) se a distância entre os termos da sequência e $x$ tender para zero. Formalmente:
  \[
    \forall \varepsilon > 0, \, \exists N\in\mathbb{N} : n\ge N \implies d(x_n,x) < \varepsilon.
  \]
\end{defbox}

\begin{defbox}[title={Continuidade em Espaços Métricos}]
  Uma função $f:(X,d_X)\to (Y,d_Y)$ é \textbf{contínua} num ponto $x_0$ se a convergência no domínio implicar a convergência no contradomínio:
  \[
    \lim_{x\to x_0} d_X(x,x_0)=0 \quad \implies \quad \lim_{x\to x_0} d_Y(f(x),f(x_0)) = 0.
  \]
  A formulação clássica equivalente (definição $\varepsilon$--$\delta$) é:
  \[
    \forall \varepsilon>0,\ \exists \delta>0: d_X(x,x_0)<\delta \implies d_Y(f(x), f(x_0))<\varepsilon.
  \]
\end{defbox}

\begin{notebox}[title={Métricas Importantes}]
  Dependendo do problema matemático, a forma de medir a distância altera-se. Exemplos notáveis:
  \begin{itemize}
      \item \textbf{Métrica Usual (Euclidiana) em $\mathbb{R}^n$:}  
      \[
        d_2(x,y)=\|x-y\|_2 = \sqrt{(x_1-y_1)^2+\cdots+(x_n-y_n)^2}
      \]
      \item \textbf{Métrica do Supremo em Funções Contínuas $C[a,b]$:}  
      \[
        d_\infty(f,g)=\sup_{x\in[a,b]} |f(x)-g(x)|
      \]
      \item \textbf{Métrica Discreta:}  
      \[
        d(x,y)=
        \begin{cases}
        0 & \text{se } x=y,\\
        1 & \text{se } x\neq y.
        \end{cases}
      \]
  \end{itemize}
\end{notebox}

% ────────────────────────────────────────────────────────
\subsection{Exemplos Resolvidos}

\begin{exbox}{1 --- Convergência na Métrica Usual}
  Considere o espaço $(\mathbb{R}, d)$ com a métrica usual $d(x,y)=|x-y|$. Mostre que a sequência definida por $x_n = \dfrac{1}{n}$ converge para $0$.
\end{exbox}
\begin{solbox}
  Queremos provar que, para qualquer $\varepsilon>0$, existe um índice $N$ a partir do qual $d(x_n, 0) < \varepsilon$.
  
  Seja $\varepsilon>0$. Como $x_n=\dfrac{1}{n}$, a distância a $0$ é:
  \[
    d(x_n,0) = \left|\frac{1}{n}-0\right| = \frac{1}{n}
  \]
  Para que $\dfrac{1}{n} < \varepsilon$, basta que $n > \dfrac{1}{\varepsilon}$. 
  
  Escolhendo um número natural $N > \dfrac{1}{\varepsilon}$ (Garantido pela Propriedade Arquimediana), temos que para qualquer $n \ge N$:
  \[
    d(x_n,0) = \frac{1}{n} \le \frac{1}{N} < \varepsilon.
  \]
  Fica assim provado que $x_n \to 0$. \checkmark
\end{solbox}

\begin{exbox}{2 --- Continuidade Extrema na Métrica Discreta}
  Mostre que \textbf{qualquer} função $f:X\to Y$ é contínua se o domínio $X$ estiver equipado com a métrica discreta.
\end{exbox}
\begin{solbox}
  Lembre-se da definição $\varepsilon$--$\delta$. Para qualquer ponto $x_0\in X$ e qualquer $\varepsilon > 0$, precisamos de encontrar um $\delta > 0$.
  
  Como $X$ tem a métrica discreta, as distâncias só podem ser $0$ ou $1$. Vamos escolher $\delta = \dfrac{1}{2}$. 
  
  Seja $x \in X$ tal que $d(x, x_0) < \dfrac{1}{2}$. Como na métrica discreta as distâncias são inteiras ($0$ ou $1$), a única forma de a distância ser menor que $0{,}5$ é se for exatamente $0$. Logo, $x = x_0$.
  
  Substituindo na função:
  \[
    d_Y(f(x), f(x_0)) = d_Y(f(x_0), f(x_0)) = 0 < \varepsilon.
  \]
  A condição verifica-se independentemente do comportamento de $f$ ou da métrica de $Y$. Assim, $f$ é sempre contínua! \checkmark
\end{solbox}

% ────────────────────────────────────────────────────────
\subsection{Exercícios Práticos}

\begin{exercisebox}{1 --- A Métrica da "Geometria do Táxi"}
  Mostre que a função $d_1(x,y)=\sum_{i=1}^n |x_i - y_i|$ define uma métrica válida em $\mathbb{R}^n$.
\end{exercisebox}
\begin{solbox}
  Temos de verificar os três axiomas:
  \begin{enumerate}
      \item \textbf{Positividade:} Como o módulo é não negativo, a soma de módulos é $\ge 0$. Além disso, $\sum |x_i - y_i| = 0 \iff \forall i, |x_i - y_i| = 0 \iff x_i = y_i \iff x = y$.
      \item \textbf{Simetria:} $|x_i - y_i| = |y_i - x_i|$, logo a soma mantém-se igual, $d_1(x,y) = d_1(y,x)$.
      \item \textbf{Desigualdade Triangular:} Usando a desigualdade triangular dos números reais ($|a+b| \le |a| + |b|$):
      \[
        |x_i - z_i| = |(x_i - y_i) + (y_i - z_i)| \le |x_i - y_i| + |y_i - z_i|
      \]
      Somando esta desigualdade de $i=1$ a $n$, obtemos imediatamente $d_1(x,z) \le d_1(x,y) + d_1(y,z)$. \checkmark
  \end{enumerate}
\end{solbox}

\begin{exercisebox}{2 --- Convergência em Espaços Discretos}
  No espaço métrico com a métrica discreta, determine todas as sequências que são convergentes.
\end{exercisebox}
\begin{solbox}
  Na métrica discreta ($d(x,y)=1$ se $x \neq y$), para uma sequência $(x_n)$ convergir para um ponto $x$, deve satisfazer: para qualquer $\varepsilon > 0$, $\exists N$ tal que $n \ge N \implies d(x_n, x) < \varepsilon$.
  
  Tomemos $\varepsilon = \dfrac{1}{2}$. Para $n \ge N$, teremos $d(x_n, x) < 0{,}5$. Pela natureza desta métrica, isto obriga a que $d(x_n, x) = 0$, ou seja, $x_n = x$.
  
  Logo, as \textbf{únicas} sequências convergentes na métrica discreta são as \textbf{sequências eventualmente constantes} (aquelas que, a partir de um certo ponto, fixam-se permanentemente no mesmo valor).
\end{solbox}

\begin{exercisebox}{3 --- Prova de Continuidade ($\varepsilon-\delta$)}
  Seja $f:\mathbb{R}\to\mathbb{R}$ dada por $f(x)=|x|$. Mostre que $f$ é contínua em todos os pontos usando a definição $\varepsilon$--$\delta$.
\end{exercisebox}
\begin{solbox}
  Pretendemos mostrar que dado $x_0 \in \mathbb{R}$ e $\varepsilon > 0$, existe $\delta > 0$ tal que $|x - x_0| < \delta \implies |f(x) - f(x_0)| < \varepsilon$.
  
  Avaliando a expressão alvo e utilizando a \textit{desigualdade triangular inversa}:
  \[
    |f(x) - f(x_0)| = |\,|x| - |x_0|\,| \le |x - x_0|
  \]
  Como o nosso objetivo é tornar a expressão menor que $\varepsilon$, basta escolhermos \textbf{$\delta = \varepsilon$}.
  Assim, se $|x - x_0| < \delta$, garantimos automaticamente que $|f(x) - f(x_0)| \le |x - x_0| < \delta = \varepsilon$. A função é contínua em $\mathbb{R}$.
\end{solbox}

\begin{exercisebox}{4 --- Limite Uniforme vs. Limite Pontual}
  Considere o espaço das funções contínuas $(C[0,1], d_\infty)$. Determine se a sequência de funções $f_n(x)=x^n$ converge uniformemente para uma função limite contínua neste espaço métrico.
\end{exercisebox}
\begin{solbox}
  Primeiro, encontramos o limite pontual de $f_n(x)$ quando $n \to \infty$:
  \[
    f(x) = \lim_{n\to\infty} x^n = 
    \begin{cases} 
      0 & \text{se } 0 \le x < 1 \\ 
      1 & \text{se } x = 1 
    \end{cases}
  \]
  Para que a sequência convirja no espaço $(C[0,1], d_\infty)$, o limite $f(x)$ \textbf{teria de pertencer} ao espaço (ou seja, teria de ser contínuo), e a distância $d_\infty(f_n, f)$ teria de ir para $0$.
  
  No entanto, a função limite $f(x)$ é descontínua em $x=1$. Mais ainda, a distância na métrica do supremo é:
  \[
    d_\infty(f_n, f) = \sup_{x \in [0,1)} |x^n - 0| = 1 \neq 0
  \]
  \textbf{Conclusão:} A sequência $f_n(x) = x^n$ \textbf{não} converge em $(C[0,1], d_\infty)$, mostrando a diferença fulcral entre convergência pontual e convergência uniforme (em métrica).
\end{solbox}

% ────────────────────────────────────────────────────────
\subsection{Questões Propostas para Defesa Oral}

\begin{oralbox}
  \textbf{Q1.} Explique a diferença entre \emph{convergência pontual} de funções e \emph{convergência em métrica} (convergência uniforme) no espaço $C[a,b]$ com a métrica do supremo.
\end{oralbox}

\begin{oralbox}
  \textbf{Q2.} Dê três exemplos de métricas diferentes em $\mathbb{R}^n$ (como a Euclideana, a da "geometria do táxi" e do Máximo) e compare geometricamente como é o aspeto das suas bolas abertas.
\end{oralbox}

\begin{oralbox}
  \textbf{Q3.} Com base na definição $\varepsilon-\delta$, explique por que razão a escolha da métrica discreta no domínio torna todas as funções matemáticas imediatamente contínuas.
\end{oralbox}

\begin{oralbox}
  \textbf{Q4.} Descreva geometricamente a desigualdade triangular. Porque é que ela é um axioma estritamente necessário para que uma função possa ser chamada de "distância" credível?
\end{oralbox}

\begin{oralbox}
  \textbf{Q5.} Num espaço métrico rigoroso, explique por que razão a convergência de uma sequência determina de forma estrita e \emph{única} o seu limite.
\end{oralbox}

\begin{oralbox}
  \textbf{Q6.} Apresente um exemplo clássico de uma sequência num espaço métrico (como os números racionais com a métrica usual) que, embora os seus termos se aproximem uns dos outros, \emph{não} converge dentro desse próprio espaço, justificando o porquê.
\end{oralbox}

\begin{oralbox}
  \textbf{Q7.} Explique a relação profunda entre \emph{bolas abertas} e o conceito topológico de \emph{continuidade}, mostrando como a pré-imagem de uma bola aberta atua na vizinhança do domínio.
\end{oralbox}
% ════════════════════════════════════════════════════════════════
% 1.4 - INTRODUÇÃO À TOPOLOGIA
% ════════════════════════════════════════════════════════════════

\section{Introdução à Topologia: Espaços, Vizinhanças e Continuidade}

% ────────────────────────────────────────────────────────
\subsection{Enquadramento Teórico}

É perfeitamente normal achar a topologia um pouco abstrata no início. Na matemática escolar, estamos muito habituados à geometria tradicional, onde medimos distâncias exatas, ângulos e comprimentos. A topologia, por outro lado, é frequentemente chamada de \textbf{``geometria da folha de borracha''}.

Ela não se importa com a distância exata entre dois pontos, mas sim com a forma como o espaço está \emph{conectado}. Num espaço topológico, podemos esticar, encolher ou torcer o espaço à vontade. A única coisa que não podemos fazer é \textbf{rasgar} ou \textbf{colar}.

\subsubsection*{Espaços Topológicos e Conjuntos Abertos}

Para falar de proximidade sem usar uma ``fita métrica'', os matemáticos inventaram o conceito de \textbf{conjunto aberto}.

\begin{defbox}[title={Espaço Topológico}]
  Um \textbf{espaço topológico} é um par $(X, \mathcal{T})$, onde:
  \begin{itemize}
    \item $X$ é um conjunto de ``coisas'' (podem ser números, pontos, pessoas, etc.).
    \item $\mathcal{T}$ (tau) é uma coleção de subconjuntos de $X$ a que chamamos \textbf{conjuntos abertos}.
  \end{itemize}
\end{defbox}

Para que $\mathcal{T}$ seja uma topologia válida, tem de obedecer a três regras fundamentais:
\begin{enumerate}
  \item \textbf{O todo e o nada estão incluídos:} O conjunto vazio $\emptyset$ e o conjunto total $X$ têm de ser considerados abertos (pertencem a $\mathcal{T}$).
  \item \textbf{Interseções finitas:} A interseção de um número \emph{finito} de conjuntos abertos é um conjunto aberto.
  \item \textbf{Uniões arbitrárias:} A união de \emph{qualquer quantidade} de conjuntos abertos é um conjunto aberto.
\end{enumerate}

\begin{notebox}[title={Intuição: Conjunto Aberto}]
  Pensa num ``conjunto aberto'' como uma zona onde, se estiveres lá dentro, podes dar um pequeno passo em qualquer direção e continuas dentro dessa zona. Não há uma fronteira rígida a limitar-te.
\end{notebox}

\subsubsection*{Vizinhanças}

A vizinhança é a forma topológica de dizer ``perto de''.

\begin{defbox}[title={Vizinhança}]
  Uma \textbf{vizinhança} de um ponto $x$ é qualquer conjunto que contenha um conjunto aberto que, por sua vez, contém o ponto $x$.
\end{defbox}

\begin{notebox}[title={Intuição: Vizinhança}]
  Imagina que o ponto $x$ és tu. Um conjunto aberto que te contém é a tua casa. Uma vizinhança tua pode ser a tua casa, mas também pode ser a tua casa mais o passeio da rua, ou até o teu bairro inteiro. Desde que o conjunto inclua a tua ``zona segura de conforto'' (o conjunto aberto), ele é considerado a tua vizinhança.
\end{notebox}

\subsubsection*{Continuidade}

No cálculo básico, aprendemos que uma função contínua é aquela que se pode desenhar ``sem levantar a caneta do papel''. Na topologia, como não temos gráficos tradicionais nem distâncias, a definição baseia-se inteiramente em conjuntos abertos.

\begin{defbox}[title={Função Contínua}]
  Uma função $f: X \to Y$ é \textbf{contínua} se, para qualquer conjunto aberto $V \subseteq Y$, a sua pré-imagem
  \[
    f^{-1}(V) = \{ x \in X \mid f(x) \in V \}
  \]
  é um conjunto aberto em $X$.
\end{defbox}

\begin{notebox}[title={Intuição: Continuidade}]
  Se pegares num grupo de pontos ``bem comportados'' (um conjunto aberto) no destino $Y$, e olhares para de onde eles vieram no espaço original $X$, eles também têm de formar um grupo bem comportado (um conjunto aberto). Isto garante que pontos que estavam ``próximos'' antes da transformação continuam estruturalmente coerentes depois. Não houve ``rasgões''.
\end{notebox}

\subsubsection*{O Donut e a Caneca de Café}

\begin{center}
  \textit{``Um topólogo é alguém que não consegue distinguir um donut da sua caneca de café.''}
\end{center}

A topologia exige uma mudança de mentalidade: da \emph{medição exata} para a \emph{relação e estrutura}. A \textbf{Regra de Ouro} dita que podemos moldar, amassar, esticar e encolher um objeto à vontade. As únicas ações estritamente proibidas são \textbf{fazer buracos novos} (rasgar) e \textbf{unir partes separadas} (colar).

\textbf{A Transformação: De Donut a Caneca} \\
Imagina que tens um donut feito de plasticina infinitamente elástica:
\begin{enumerate}
  \item \textbf{Ponto de partida:} Tens o donut, que tem exatamente \emph{um buraco} no meio.
  \item \textbf{Criar o copo:} Beliscas e empurras uma das metades da massa, criando uma grande depressão. Como não atravessas a massa para o outro lado, não estás a fazer um buraco novo — apenas a esticar a superfície.
  \item \textbf{Formar a pega:} A parte restante do donut (com o buraco original intacto) é afinada e esticada, tornando-se a pega da caneca.
\end{enumerate}
O buraco no meio do donut tornou-se no buraco por onde passas os dedos na pega da caneca.

\begin{defbox}[title={Homeomorfismo}]
  Dois espaços são \textbf{topologicamente equivalentes} (homeomorfos) se for possível transformar um no outro sem rasgar nem colar. Esta equivalência (uma função bijetiva e contínua com inversa contínua) chama-se \textbf{homeomorfismo}.
\end{defbox}

O que define a ``forma'' de um espaço em topologia não é o volume ou os ângulos, mas os seus \textbf{invariantes topológicos} — neste caso, o \emph{número de buracos} (chamado \textbf{género}).

\begin{center}
  \renewcommand{\arraystretch}{1.4}
  \begin{tabular}{lcl}
    \toprule
    \textbf{Objeto} & \textbf{Género (n.º de buracos)} & \textbf{Equivalente a} \\
    \midrule
    Donut (toro) & 1 & Caneca de café \\
    Caneca de café & 1 & Donut (toro) \\
    Bola de bilhar & 0 & Esfera, cubo, prato \ldots \\
    \bottomrule
  \end{tabular}
\end{center}

\begin{notebox}[title={Intuição sobre Buracos}]
  O espaço interior da caneca onde entra o líquido \emph{não} é um buraco topológico — é apenas uma reentrância cega (uma ``mossa''). O único buraco topológico da caneca é o da pega.
\end{notebox}

% ────────────────────────────────────────────────────────
\subsection{Exemplos Resolvidos}

\begin{exbox}{1 --- Topologia Finita}
  Imagina um conjunto com três pessoas: $X = \{\text{Ana, Bruno, Carlos}\}$. Define-se a coleção:
  \[
    \mathcal{T} = \bigl\{\, \emptyset,\; \{\text{Ana, Bruno, Carlos}\},\; \{\text{Ana}\} \,\bigr\}
  \]
  Verifique se $\mathcal{T}$ é uma topologia e se $\{\text{Ana, Bruno}\}$ é uma vizinhança da Ana.
\end{exbox}
\begin{solbox}
  \textbf{1. Verificação da Topologia:}
  \begin{itemize}
    \item $\emptyset$ e $X$ pertencem a $\mathcal{T}$. \checkmark
    \item A interseção de $\{\text{Ana}\}$ com $X$ é $\{\text{Ana}\}$, que está em $\mathcal{T}$. \checkmark
    \item A união de $\{\text{Ana}\}$ com $\emptyset$ é $\{\text{Ana}\}$, que está em $\mathcal{T}$. \checkmark
  \end{itemize}
  Logo, $(X, \mathcal{T})$ é um espaço topológico válido.
  
  \textbf{2. Vizinhança:}
  O conjunto $\{\text{Ana, Bruno}\}$ é uma \emph{vizinhança} da Ana? \textbf{Sim}, porque contém o conjunto aberto $\{\text{Ana}\}$, e esse conjunto aberto contém a Ana.
\end{solbox}

\begin{exbox}{2 --- A Reta Real (Topologia Padrão)}
  Na reta real $\mathbb{R}$, os conjuntos abertos são os intervalos abertos (sem as extremidades). Analise a continuidade da função $f(x) = 2x$.
\end{exbox}
\begin{solbox}
  Na topologia padrão, o intervalo $(0, 1)$ — que inclui $0{,}1$, $0{,}5$, $0{,}999$, mas \textbf{não} o $0$ nem o $1$ — é um conjunto aberto. Para qualquer ponto interior (ex: $0{,}99$), existe sempre um intervalo menor em seu redor que o contém por completo e ainda está dentro de $(0,1)$.

  A função $f(x) = 2x$ é contínua. Para provar topologicamente, basta ver que a pré-imagem de um conjunto aberto qualquer, como o intervalo $V = (2, 4)$ no destino $Y$, é o conjunto $f^{-1}(V) = (1, 2)$ no domínio $X$. Como $(1, 2)$ também é um intervalo aberto, a pré-imagem de um aberto é aberta. A estrutura mantém-se!
\end{solbox}

% ────────────────────────────────────────────────────────
\subsection{Exercícios Práticos}

\begin{exercisebox}{1}
  Seja $X = \{1, 2, 3\}$. Considere a coleção $\mathcal{T} = \{\emptyset, \{1\}, \{2\}, \{1, 2, 3\}\}$. Verifique se $\mathcal{T}$ forma uma topologia em $X$. Em caso negativo, explique qual a regra que falha.
\end{exercisebox}
\begin{solbox}
  Para ser topologia, $\mathcal{T}$ tem de ser fechada para uniões arbitrárias. 
  
  Se unirmos o conjunto aberto $\{1\}$ com o conjunto aberto $\{2\}$, obtemos:
  \[ \{1\} \cup \{2\} = \{1, 2\} \]
  
  No entanto, o conjunto $\{1, 2\}$ \textbf{não pertence} à coleção $\mathcal{T}$. Como a regra da união falha, $\mathcal{T}$ \textbf{não} é uma topologia válida em $X$.
\end{solbox}

% ────────────────────────────────────────────────────────
\subsection{Questões Propostas para Defesa Oral}

\begin{oralbox}
  \textbf{Q1.} Explique, por palavras suas e usando a analogia da "geometria da folha de borracha", por que motivo a topologia não utiliza conceitos como "distância" ou "ângulo" e como o conceito de \emph{conjunto aberto} substitui a fita métrica para definir a proximidade.
\end{oralbox}

\begin{oralbox}
  \textbf{Q2.} Um prato de sopa fundo e uma bola de bowling são topologicamente equivalentes (homeomorfos)? Justifique a sua resposta com base no conceito de género (número de buracos) e nas regras de deformação topológica permitidas.
\end{oralbox}
% ════════════════════════════════════════════════════════════════
% 0.5 - LIMITES E CONTINUIDADE EM ESPAÇOS MÉTRICOS E TOPOLÓGICOS
% ════════════════════════════════════════════════════════════════

\section{Limites e Continuidade de Funções}

% ────────────────────────────────────────────────────────
\subsection{Enquadramento Teórico}

\begin{notebox}[title={Intuição Geral}]
  O conceito de \textbf{limite} descreve a ideia fundamental de que, quando os pontos do domínio de uma função se aproximam cada vez mais de um certo valor, as imagens desses pontos também se aproximam de um valor determinado.

  A \textbf{continuidade} não é mais do que a formalização matemática rigorosa da seguinte frase:
  \begin{center}
      \textit{``Pequenas variações na entrada produzem apenas pequenas variações na saída.''}
  \end{center}
\end{notebox}

\subsubsection*{Definição de Limite em Espaços Métricos}

Sejam $(X,d_X)$ e $(Y,d_Y)$ espaços métricos, com uma função $f : X \to Y$, e seja $a \in X$ um ponto de acumulação.

\begin{defbox}[title={Limite de uma Função ($\varepsilon-\delta$)}]
  Dizemos que o \textbf{limite de $f(x)$ quando $x$ tende para $a$ é $L$}, denotado por $\lim_{x \to a} f(x) = L$, se para qualquer tolerância de erro no contradomínio ($\varepsilon$), existir uma margem de aproximação no domínio ($\delta$) tal que:
  \[
    \forall \varepsilon > 0, \;\exists \delta > 0 : \quad 0 < d_X(x,a) < \delta \implies d_Y(f(x),L) < \varepsilon
  \]
\end{defbox}

\begin{notebox}[title={Como Interpretar}]
  A definição diz-nos que, controlando \emph{o quão perto} $x$ está de $a$ (através do raio $\delta$), conseguimos forçar e controlar \emph{o quão perto} $f(x)$ está do limite $L$ (dentro do erro $\varepsilon$).
\end{notebox}

\subsubsection*{Continuidade em Espaços Métricos e Topológicos}

\begin{defbox}[title={Continuidade Num Ponto}]
  Uma função $f:X\to Y$ é \textbf{contínua no ponto $a$} se o limite da função nesse ponto coincidir exatamente com o valor da função no ponto:
  \[
    \lim_{x\to a} f(x)=f(a)
  \]
  Substituindo na definição $\varepsilon$--$\delta$, temos:
  \[
    \forall \varepsilon>0, \;\exists \delta>0 : \quad d_X(x,a)<\delta \implies d_Y(f(x),f(a))<\varepsilon
  \]
\end{defbox}

\begin{thmbox}[title={Critério Topológico Global}]
  Se passarmos de espaços métricos para espaços topológicos mais abstratos $(X,\tau_X)$ e $(Y,\tau_Y)$, a métrica de distância desaparece. 
  Neste caso, $f:X\to Y$ é contínua se, e só se, a \textbf{imagem inversa de qualquer conjunto aberto for um conjunto aberto}:
  \[
    f^{-1}(V) \text{ é aberto em } X, \quad \forall V \subseteq Y \text{ aberto.}
  \]
\end{thmbox}

% ────────────────────────────────────────────────────────
\subsection{Exemplos Resolvidos}

\begin{exbox}{1 --- Continuidade de uma Função Afim}
  Considere a função $f:\mathbb{R}\to\mathbb{R}$ dada por $f(x)=3x+1$. Mostrar pela definição $\varepsilon-\delta$ que $f$ é contínua em todo o $\mathbb{R}$.
\end{exbox}
\begin{solbox}
  Seja $a \in \mathbb{R}$ um ponto qualquer e $\varepsilon > 0$. Queremos encontrar $\delta > 0$ tal que $|x-a| < \delta \implies |f(x)-f(a)| < \varepsilon$.
  
  Analisamos a distância no contradomínio:
  \[
    |f(x)-f(a)| = |(3x+1) - (3a+1)| = |3x - 3a| = 3|x-a|
  \]
  Como queremos que $3|x-a| < \varepsilon$, basta isolar a distância do domínio, o que nos sugere que $|x-a| < \frac{\varepsilon}{3}$.
  
  Portanto, basta tomar \textbf{$\delta = \frac{\varepsilon}{3}$}. Assim, se $|x-a| < \delta$, temos garantidamente:
  \[
    |f(x)-f(a)| = 3|x-a| < 3\left(\frac{\varepsilon}{3}\right) = \varepsilon
  \]
  Logo, $f$ é contínua em qualquer $a \in \mathbb{R}$. \checkmark
\end{solbox}

\begin{exbox}{2 --- Limite de uma Função Raiz}
  Mostrar formalmente que $\lim_{x\to 2} \sqrt{x} = \sqrt{2}$.
\end{exbox}
\begin{solbox}
  Queremos provar que $\forall \varepsilon > 0, \exists \delta > 0 : |x-2| < \delta \implies |\sqrt{x}-\sqrt{2}| < \varepsilon$.
  
  Começamos por manipular algebricamente a expressão do contradomínio multiplicando pelo conjugado:
  \[
    |\sqrt{x}-\sqrt{2}| = \left| \frac{(\sqrt{x}-\sqrt{2})(\sqrt{x}+\sqrt{2})}{\sqrt{x}+\sqrt{2}} \right| = \frac{|x-2|}{\sqrt{x}+\sqrt{2}}
  \]
  Para isolar $|x-2|$, precisamos de minorar o denominador. Se restringirmos $x$ suficientemente perto de $2$, digamos impondo $|x-2| < 1$, então $x \in (1, 3)$. 
  Isto implica que $\sqrt{x} > \sqrt{1} = 1$, logo:
  \[
    \sqrt{x}+\sqrt{2} > 1+\sqrt{2}
  \]
  Substituindo na nossa fração, obtemos:
  \[
    |\sqrt{x}-\sqrt{2}| = \frac{|x-2|}{\sqrt{x}+\sqrt{2}} < \frac{|x-2|}{1+\sqrt{2}}
  \]
  Para garantir que a expressão inteira seja menor que $\varepsilon$, basta tomar:
  \[
    \delta = \min\bigl\{1, \, (1+\sqrt{2})\varepsilon\bigr\}
  \]
  O limite está provado. \checkmark
\end{solbox}

% ────────────────────────────────────────────────────────
\subsection{Exercícios Práticos}

\begin{exercisebox}{1 --- Verificação com $\varepsilon-\delta$ para Quadráticas}
  Verifique usando a definição formal $\varepsilon$--$\delta$ que $\lim_{x\to 1}(2x^2+3)=5$.
\end{exercisebox}
\begin{solbox}
  Queremos que $| (2x^2+3) - 5 | < \varepsilon$, o que simplifica para $2|x^2-1| < \varepsilon \iff 2|x-1||x+1| < \varepsilon$.
  
  Para limitar o termo $|x+1|$, impomos provisoriamente $\delta \le 1$.
  Se $|x-1| < 1 \implies 0 < x < 2 \implies 1 < x+1 < 3 \implies |x+1| < 3$.
  
  Substituindo esta limitação na nossa equação inicial:
  \[
    2|x-1||x+1| < 2|x-1|(3) = 6|x-1|
  \]
  Queremos que $6|x-1| < \varepsilon$, ou seja, $|x-1| < \frac{\varepsilon}{6}$.
  
  Basta escolher $\delta = \min\left\{1, \frac{\varepsilon}{6}\right\}$. Assim, o limite está provado formalmente.
\end{solbox}

\begin{exercisebox}{2 --- Continuidade do Módulo}
  Mostre que a função $f(x)=|x|$ é contínua em todo o $\mathbb{R}$.
\end{exercisebox}
\begin{solbox}
  Queremos mostrar que $\forall a \in \mathbb{R}$, $\lim_{x\to a} |x| = |a|$.
  
  Pela propriedade da Desigualdade Triangular Inversa, sabemos que para quaisquer reais $x$ e $a$:
  \[
    \bigl| |x| - |a| \bigr| \le |x-a|
  \]
  Dado um $\varepsilon > 0$, basta escolher $\delta = \varepsilon$. 
  Deste modo, se $|x-a| < \delta$, então:
  \[
    |f(x) - f(a)| = \bigl| |x| - |a| \bigr| \le |x-a| < \delta = \varepsilon
  \]
  Isto prova a continuidade de $f(x)=|x|$ para qualquer ponto do domínio.
\end{solbox}

\begin{exercisebox}{3 --- Continuidade de Funções por Ramos}
  Determine se a função $f:\mathbb{R}\to\mathbb{R}$ definida por
  \[
    f(x)=\begin{cases}
    x^2, & x\neq 1\\
    3, & x=1
    \end{cases}
  \]
  é contínua em $x=1$. Justifique a resposta.
\end{exercisebox}
\begin{solbox}
  Para ser contínua em $x=1$, é estritamente necessário que $\lim_{x \to 1} f(x) = f(1)$.
  
  Calculando o limite (que analisa os valores perto de $x=1$, mas não no próprio ponto $x=1$):
  \[
    \lim_{x \to 1} f(x) = \lim_{x \to 1} x^2 = 1^2 = 1
  \]
  Contudo, a definição da função dita que:
  \[
    f(1) = 3
  \]
  Como $\lim_{x \to 1} f(x) \neq f(1)$ (uma vez que $1 \neq 3$), concluímos que a função \textbf{não é contínua} em $x=1$. Tem uma descontinuidade removível.
\end{solbox}

\begin{exercisebox}{4 --- A Curva de Seno do Topólogo}
  Considere a função $f:\mathbb{R} \to \mathbb{R}$ dada por $f(x)=\sin(1/x)$ para $x\neq 0$, com $f(0)=0$. A função é contínua em $x=0$? Justifique.
\end{exercisebox}
\begin{solbox}
  Para ser contínua em $x=0$, teria de verificar-se $\lim_{x \to 0} \sin(1/x) = 0$.
  
  No entanto, à medida que $x \to 0$, o argumento $1/x \to \pm \infty$. A função seno irá oscilar infinitamente entre $-1$ e $1$, cada vez mais rápido, sem nunca se fixar ou convergir para um único valor.
  
  Prova rápida por sucessões (Critério de Heine):
  \begin{itemize}
      \item Se escolhermos $x_n = \frac{1}{\pi/2 + 2n\pi} \to 0$, $f(x_n) = \sin(\pi/2 + 2n\pi) = 1 \to 1$.
      \item Se escolhermos $y_n = \frac{1}{n\pi} \to 0$, $f(y_n) = \sin(n\pi) = 0 \to 0$.
  \end{itemize}
  Como trajetórias diferentes geram limites diferentes ($1 \neq 0$), o limite $\lim_{x \to 0} f(x)$ não existe. Logo, a função \textbf{não é contínua} em $x=0$.
\end{solbox}

% ────────────────────────────────────────────────────────
\subsection{Questões Propostas para Defesa Oral}

\begin{oralbox}
  \textbf{Q1.} Explique, usando a analogia de jogo ou tolerâncias, o que significa intuitivamente a definição formal $\varepsilon-\delta$ do limite de uma função num ponto.
\end{oralbox}

\begin{oralbox}
  \textbf{Q2.} Em que sentido exato é que a \emph{continuidade} é apenas uma versão restrita e ``sem salto'' do conceito geral de limite?
\end{oralbox}

\begin{oralbox}
  \textbf{Q3.} Justifique o porquê da definição topológica de continuidade (a pré-imagem de um conjunto aberto no contradomínio ter de ser necessariamente um conjunto aberto no domínio) ser matematicamente equivalente à definição $\varepsilon$--$\delta$ clássica quando aplicada em espaços métricos.
\end{oralbox}

\begin{oralbox}
  \textbf{Q4.} Dê um exemplo prático ou geométrico de uma função descontínua e explique a razão lógica / topológica de ocorrência dessa quebra de continuidade.
\end{oralbox}

\begin{oralbox}
  \textbf{Q5.} Qual a relação entre a continuidade global de funções e a preservação do limite na convergência de sucessões (Critério de Heine / Teorema de Continuidade por Sucessões)?
\end{oralbox}
% ════════════════════════════════════════════════════════════════
% SECÇÃO - PERSPETIVA ESTRUTURAL
% ════════════════════════════════════════════════════════════════

\section{Perspetiva Estrutural}

% ────────────────────────────────────────────────────────
\subsection{Enquadramento Teórico}

A perspetiva estrutural da matemática, frequentemente formalizada através da Teoria das Categorias, propõe uma mudança de paradigma: em vez de olharmos para os elementos internos que compõem uma entidade matemática (como os números dentro de um conjunto), focamo-nos nas relações e nas transformações entre essas entidades.

\subsubsection*{1. Objetos (Objects)}

Na teoria das categorias, os objetos são os blocos de construção primários. Do ponto de vista estrutural, não nos interessa o que está ``dentro'' do objeto, mas sim como ele interage com os outros.

\begin{defbox}[title={Objeto}]
  Um \textbf{objeto} é uma entidade abstrata que pertence a uma determinada ``categoria''.
\end{defbox}

\textbf{Exemplos de Objetos em Categorias:}
\begin{itemize}
  \item Na categoria dos Conjuntos (\textbf{Set}), os objetos são conjuntos (ex: $\{1, 2, 3\}$, $\mathbb{N}$, $\mathbb{R}$).
  \item Na categoria dos Espaços Vetoriais (\textbf{Vect}), os objetos são espaços vetoriais (ex: $\mathbb{R}^2$, $\mathbb{R}^3$).
  \item Na categoria dos Espaços Topológicos (\textbf{Top}), os objetos são espaços topológicos.
\end{itemize}

\begin{notebox}[title={Interpretação Visual}]
  Geometricamente ou num diagrama, um objeto é representado simplesmente como um ponto (ou um nó).
\end{notebox}

\subsubsection*{2. Morfismos (Mappings)}

Se os objetos são as entidades, os morfismos são os ``caminhos'' ou transformações estruturais entre eles.

\begin{defbox}[title={Morfismo}]
  Um \textbf{morfismo} $f$ é uma relação direcional entre um objeto de origem $A$ (domínio) e um objeto de destino $B$ (contradomínio), denotado por $f: A \to B$. Em muitas categorias, estes morfismos são funções que preservam a estrutura matemática dos objetos.
\end{defbox}

\textbf{Exemplos de Morfismos em Categorias:}
\begin{itemize}
  \item Em \textbf{Set}, os morfismos são funções normais entre conjuntos.
  \item Em \textbf{Vect}, os morfismos são transformações lineares (preservam a adição de vetores e a multiplicação por escalares).
  \item Em \textbf{Top}, os morfismos são funções contínuas (preservam a estrutura de vizinhança/limites, fundamental em Análise Matemática).
\end{itemize}

\begin{notebox}[title={Interpretação Visual}]
  Os morfismos são desenhados como setas direcionadas que ligam um ponto (objeto) a outro.
\end{notebox}

\subsubsection*{3. Composição}

A verdadeira essência da perspetiva estrutural surge na composição: a capacidade de encadear aplicações de forma consistente.

\begin{defbox}[title={Composição e Axiomas da Categoria}]
  Dados três objetos $A$, $B$ e $C$, e dois morfismos $f: A \to B$ e $g: B \to C$, existe um morfismo composto único, denotado por $g \circ f: A \to C$ (lê-se ``$g$ após $f$''). 
  
  Para que uma estrutura seja uma categoria válida, a composição tem de obedecer a duas regras:
  \begin{enumerate}
      \item \textbf{Associatividade:} $h \circ (g \circ f) = (h \circ g) \circ f$.
      \item \textbf{Identidade:} Para cada objeto $X$, existe um morfismo de identidade $\text{id}_X: X \to X$ tal que, para qualquer $f: A \to B$, temos $f \circ \text{id}_A = f$ e $\text{id}_B \circ f = f$.
  \end{enumerate}
\end{defbox}

\textbf{Exemplo Prático:} \\
Se $A, B, C$ são subconjuntos de $\mathbb{R}$, e $f(x) = x^2$ e $g(y) = y + 1$, então a composição $g(f(x))$ representa o morfismo $g \circ f$ de $A$ para $C$, dado por $(g \circ f)(x) = x^2 + 1$.

\begin{notebox}[title={Interpretação Visual (Diagramas Comutativos)}]
  A composição é frequentemente visualizada através de diagramas comutativos. Se tivermos uma seta de $A$ para $B$, e de $B$ para $C$, o diagrama ``comuta'' se a seta direta de $A$ para $C$ representar o mesmo ``resultado final'' que percorrer as duas setas anteriores. Forma-se, assim, um triângulo.
\end{notebox}


% ────────────────────────────────────────────────────────
\subsection{Exemplos Resolvidos}

\begin{exbox}{1 --- Categoria de Caminhos em $\mathbb{R}^2$}
  Numa categoria onde os objetos são pontos em $\mathbb{R}^2$ e os morfismos são caminhos entre esses pontos:
  \begin{itemize}
      \item O que é a identidade $\text{id}_A$?
      \item O que representa a composição $g \circ f$?
  \end{itemize}
\end{exbox}
\begin{solbox}
  \textbf{Identidade $\text{id}_A$:} É o caminho constante (permanecer estático no próprio ponto $A$). \\
  \textbf{Composição $g \circ f$:} É a concatenação dos caminhos. Colocamos o início do caminho $g$ no final do caminho $f$. O resultado final é um novo caminho direto que nos leva de $A$ até $C$, passando obrigatoriamente por $B$.
\end{solbox}

\begin{exbox}{2 --- Verificação de Associatividade}
  Verifique a propriedade da associatividade para as seguintes funções em $\mathbb{R}$:
  \[ f(x) = 2x, \quad g(x) = x+3, \quad h(x) = x^2 \]
\end{exbox}
\begin{solbox}
  Para verificar a associatividade, precisamos de provar que $[h \circ (g \circ f)](x) = [(h \circ g) \circ f](x)$.
  
  \textbf{Cálculo do lado esquerdo:}
  \begin{align*}
      (g \circ f)(x) &= g(2x) = 2x + 3 \\
      [h \circ (g \circ f)](x) &= h(2x + 3) = \mathbf{(2x + 3)^2}
  \end{align*}
  
  \textbf{Cálculo do lado direito:}
  \begin{align*}
      (h \circ g)(x) &= h(x + 3) = (x + 3)^2 \\
      [(h \circ g) \circ f](x) &= (h \circ g)(2x) = \mathbf{(2x + 3)^2}
  \end{align*}
  
  Como ambos os caminhos de composição resultam na mesma expressão $\mathbf{(2x + 3)^2}$, a propriedade associativa é verificada. \checkmark
\end{solbox}


% ────────────────────────────────────────────────────────
\subsection{Exercícios Práticos}

\begin{exercisebox}{1}
  Considere a categoria \textbf{Set} (Conjuntos). Sejam os objetos $A = \{1, 2\}$ e $B = \{a, b, c\}$. Quantos morfismos (funções) diferentes existem de $A$ para $B$? E de $B$ para $A$?
\end{exercisebox}
\begin{solbox}
  Em \textbf{Set}, os morfismos são funções.
  \begin{itemize}
      \item O número de funções de $A$ para $B$ é dado por $|B|^{|A|}$. Como $|B|=3$ e $|A|=2$, existem $3^2 = \mathbf{9}$ morfismos de $A$ para $B$.
      \item O número de funções de $B$ para $A$ é dado por $|A|^{|B|}$. Como $|A|=2$ e $|B|=3$, existem $2^3 = \mathbf{8}$ morfismos de $B$ para $A$.
  \end{itemize}
\end{solbox}


% ────────────────────────────────────────────────────────
\subsection{Questões Propostas para Defesa Oral}

\begin{oralbox}
  \textbf{Q1.} Explique, por palavras suas, qual é a principal vantagem da mudança de paradigma oferecida pela Teoria das Categorias (foco nos \emph{morfismos}) em comparação com a Teoria de Conjuntos tradicional (foco nos \emph{elementos} internos dos objetos).
\end{oralbox}

\begin{oralbox}
  \textbf{Q2.} Num diagrama comutativo triangular com os objetos $A$, $B$ e $C$, o que significa matematicamente dizer que o diagrama "comuta"? Utilize o conceito de composição de morfismos na sua resposta.
\end{oralbox}

% ---------------------------------------------------------
% CAPÍTULO 1
% ---------------------------------------------------------
% ════════════════════════════════════════════════════════════════
% 1.1 - CONCEITOS DE DERIVADA
% ════════════════════════════════════════════════════════════════

\chapter{Cálculo Diferencial em R}

\section{Conceitos de Derivada}

% ────────────────────────────────────────────────────────
\subsection{Enquadramento Teórico}

\begin{notebox}[title={Intuição Essencial}]
  A derivada mede a \textbf{velocidade instantânea de mudança}. Se uma função descreve uma quantidade que varia (como a posição, a temperatura ou o custo financeiro), a derivada indica-nos quão rápido essa quantidade está a mudar num exato instante.
\end{notebox}

\begin{defbox}[title={Derivada de uma Função num Ponto}]
  Seja $f:\mathbb{R}\to\mathbb{R}$. Diz-se que $f$ é \textbf{derivável} (ou diferenciável) num ponto $x_0$ se o seguinte limite existir e for um número finito:
  \[
    f'(x_0) = \lim_{h\to 0} \frac{f(x_0+h)-f(x_0)}{h}
  \]
  Geometricamente, este valor representa a \textbf{inclinação (ou declive) da reta tangente} à curva representativa de $f$ no ponto de abcissa $x_0$. O termo interno $\frac{f(x_0+h)-f(x_0)}{h}$ é a \emph{taxa média de variação} (ou quociente incremental).
\end{defbox}

\begin{defbox}[title={Regras Fundamentais de Derivação}]
  Sejam $f$ e $g$ funções deriváveis e $c, n \in \mathbb{R}$. As operações com derivadas respeitam as seguintes regras operatórias:
  \begin{itemize}
      \item \textbf{Constante:} $(c)' = 0$
      \item \textbf{Potência:} $(x^n)' = n x^{n-1}$
      \item \textbf{Soma e Diferença:} $(f \pm g)' = f' \pm g'$
      \item \textbf{Multiplicação por Escalar:} $(cf)' = c f'$
      \item \textbf{Produto:} $(f \cdot g)' = f'g + f g'$
      \item \textbf{Quociente:} $\displaystyle \left(\frac{f}{g}\right)' = \frac{f'g - f g'}{g^2}$ \quad (para $g \neq 0$)
      \item \textbf{Regra da Cadeia (Funções Compostas):} $(f \circ g)'(x) = f'(g(x)) \cdot g'(x)$
  \end{itemize}
\end{defbox}

\begin{warnbox}[title={Erros Comuns na Derivação}]
  \begin{itemize}
      \item \textbf{Confundir o quociente incremental com a derivada:} A expressão $\frac{f(x+h)-f(x)}{h}$ dá-nos apenas a velocidade \emph{média} entre dois pontos. A derivada (velocidade \emph{instantânea}) exige a aplicação obrigatória do limite $h \to 0$.
      \item \textbf{Achar que "Contínua $\implies$ Derivável":} Esta implicação é \textbf{falsa}. Toda a função derivável é garantidamente contínua, mas \emph{nem toda} função contínua é derivável. O contraexemplo clássico é $f(x) = |x|$, que não tem derivada na origem devido a uma "quebra" abrupta na geometria.
  \end{itemize}
\end{warnbox}

% ────────────────────────────────────────────────────────
\subsection{Exemplos Resolvidos}

\begin{exbox}{1 --- Derivada de uma Função Potência}
  Calcular a derivada da função $f(x)=x^3$ e interpretar o resultado.
\end{exbox}
\begin{solbox}
  Aplicando diretamente a regra da potência $(x^n)' = n x^{n-1}$ com $n=3$:
  \[
    f'(x) = 3x^{3-1} = \mathbf{3x^2}
  \]
  \textbf{Interpretação:} A inclinação da reta tangente à curva num dado ponto $x$ é dada por $3x^2$. Como $x^2 \ge 0$, a função está sempre a crescer (ou com inclinação nula na origem, onde $x=0$). À medida que o valor de $x$ se afasta do centro, a inclinação aumenta exponencialmente, tornando a curva cada vez mais íngreme.
\end{solbox}

\begin{exbox}{2 --- Falha na Derivabilidade (O "Bico")}
  A função $f(x)=|x|$ não é derivável em $x=0$. Explique o porquê de forma intuitiva.
\end{exbox}
\begin{solbox}
  A função módulo forma um "bico" (ponto anguloso) exatamente em $x=0$.
  Se analisarmos as inclinações das semirretas laterais:
  \begin{itemize}
      \item O lado esquerdo da curva tem inclinação constante de $-1$.
      \item O lado direito da curva tem inclinação constante de $+1$.
  \end{itemize}
  Como as inclinações laterais não coincidem ($-1 \neq 1$), o limite que define a derivada não existe nesse ponto. Na prática, não conseguimos equilibrar uma \emph{única} reta tangente na ponta daquele bico.
\end{solbox}

% ────────────────────────────────────────────────────────
\subsection{Exercícios Práticos}

\begin{exercisebox}{1 --- Aplicação Múltipla de Regras}
  Calcule a função derivada $f'(x)$ para o polinómio $f(x)=5x^4 - 3x + 7$.
\end{exercisebox}
\begin{solbox}
  Utilizando linearmente as regras da soma, da multiplicação por escalar, da potência e da constante:
  \begin{align*}
      f'(x) &= (5x^4)' - (3x)' + (7)' \\
            &= 5(4x^3) - 3(1) + 0 \\
            &= 20x^3 - 3
  \end{align*}
  \textbf{Solução Final:} $\boxed{f'(x) = 20x^3 - 3}$.
\end{solbox}

\begin{exercisebox}{2 --- Análise Analítica de Derivabilidade}
  Determine se a função $f(x)=|x-2|$ é derivável em $x=2$. Justifique o seu raciocínio.
\end{exercisebox}
\begin{solbox}
  Vamos calcular os limites laterais do quociente incremental no ponto $x_0 = 2$.
  Sabemos que $f(2) = |2-2| = 0$.
  
  \textbf{Limite à direita ($h \to 0^+$):}
  \[
    \lim_{h \to 0^+} \frac{f(2+h) - f(2)}{h} = \lim_{h \to 0^+} \frac{|2+h-2|}{h} = \lim_{h \to 0^+} \frac{|h|}{h} = \lim_{h \to 0^+} \frac{h}{h} = \mathbf{1}
  \]
  
  \textbf{Limite à esquerda ($h \to 0^-$):}
  \[
    \lim_{h \to 0^-} \frac{f(2+h) - f(2)}{h} = \lim_{h \to 0^-} \frac{|h|}{h} = \lim_{h \to 0^-} \frac{-h}{h} = \mathbf{-1}
  \]
  
  Como os limites laterais devolvem valores distintos ($1 \neq -1$), o limite global em $x=2$ não existe. \\
  \textbf{Conclusão:} A função \textbf{não é derivável} em $x=2$.
\end{solbox}

\begin{exercisebox}{3 --- Prova pela Definição Formal}
  Mostre, recorrendo estritamente à definição formal de limite, que para $f(x) = x^2$, a derivada é efetivamente $(x^2)' = 2x$.
\end{exercisebox}
\begin{solbox}
  Pela definição, a derivada é dada por:
  \[
    f'(x) = \lim_{h\to 0} \frac{f(x+h)-f(x)}{h}
  \]
  Substituindo a função pela expressão $f(x) = x^2$:
  \begin{align*}
      f'(x) &= \lim_{h\to 0} \frac{(x+h)^2 - x^2}{h} \\
            &= \lim_{h\to 0} \frac{x^2 + 2xh + h^2 - x^2}{h} \\
            &= \lim_{h\to 0} \frac{2xh + h^2}{h}
  \end{align*}
  Fatorizando o $h$ no numerador (visto que $h \neq 0$ durante a aproximação no limite):
  \[
    f'(x) = \lim_{h\to 0} \frac{h(2x + h)}{h} = \lim_{h\to 0} (2x + h)
  \]
  Aplicando a substituição direta do limite quando $h \to 0$, obtemos finalmente:
  \[
    f'(x) = 2x + 0 = \mathbf{2x} \quad \text{\checkmark}
  \]
\end{solbox}

% ────────────────────────────────────────────────────────
\subsection{Questões Propostas para Defesa Oral}

\begin{oralbox}
  \textbf{Q1.} Como explicaria o conceito prático de derivada a alguém do ensino básico? (Sugestão: recorra à analogia de um carro em viagem a medir a velocidade no conta-quilómetros).
\end{oralbox}

\begin{oralbox}
  \textbf{Q2.} Qual a diferença rigorosa entre \emph{taxa média de variação} e \emph{taxa instantânea de variação}? Relacione ambos os conceitos geométricos (reta secante \textit{vs.} reta tangente).
\end{oralbox}

\begin{oralbox}
  \textbf{Q3.} Em que pontos exatos a função modular típica, como $|x|$, falha a derivabilidade e porquê? Explique porque é que a \emph{continuidade} não é suficiente para garantir a derivação.
\end{oralbox}

\begin{oralbox}
  \textbf{Q4.} Explique geometricamente ou através de taxas de variação interdependentes a intuição por trás da Regra da Cadeia. Como se comportam os declives numa função composta?
\end{oralbox}
% ════════════════════════════════════════════════════════════════
% 1.2 - DERIVADAS DE FUNÇÕES TRIGONOMÉTRICAS INVERSAS
% ════════════════════════════════════════════════════════════════

\section{Derivadas de Funções Trigonométricas Inversas}

% ────────────────────────────────────────────────────────
\subsection{Enquadramento Teórico}

\begin{notebox}[title={Intuição Essencial}]
  As funções trigonométricas inversas ($\arcsin$, $\arccos$, $\arctan$) permitem recuperar o \textbf{ângulo} a partir da sua respetiva razão trigonométrica. As derivadas destas funções são ferramentas analíticas essenciais quando lidamos com expressões que envolvem ângulos e queremos simplificá-las para formas puramente algébricas.
\end{notebox}

\begin{defbox}[title={Derivadas Fundamentais}]
  As derivadas das principais funções trigonométricas inversas para a variável simples $x$ são dadas por:
  \begin{align*}
      \frac{\dd}{\dd x} (\arcsin x) &= \frac{1}{\sqrt{1-x^2}}, \qquad |x| < 1 \\[1ex]
      \frac{\dd}{\dd x} (\arccos x) &= -\frac{1}{\sqrt{1-x^2}}, \qquad |x| < 1 \\[1ex]
      \frac{\dd}{\dd x} (\arctan x) &= \frac{1}{1+x^2}, \qquad x \in \mathbb{R}
  \end{align*}
\end{defbox}

\begin{defbox}[title={Generalização com a Regra da Cadeia}]
  Na prática, raramente derivamos apenas $x$. Se $u(x)$ for uma função derivável em ordem a $x$, as fórmulas acima generalizam-se multiplicando pela derivada do argumento $u'(x)$:
  \begin{align*}
      (\arcsin u)' &= \frac{u'}{\sqrt{1-u^2}} \\[1ex]
      (\arccos u)' &= -\frac{u'}{\sqrt{1-u^2}} \\[1ex]
      (\arctan u)' &= \frac{u'}{1+u^2}
  \end{align*}
\end{defbox}

\begin{notebox}[title={Dicas de Memorização e Origem Geométrica}]
  \begin{itemize}
      \item $\arcsin$ e $\arccos$ partilham a \textbf{exata mesma fórmula}, mudando apenas o sinal (o $\arccos$ é negativo).
      \item O $\arctan$ é o mais simples de memorizar, pois não tem raízes quadradas: \(\frac{1}{1+x^2}\).
      \item As raízes quadradas em $\arcsin$ e $\arccos$ aparecem naturalmente devido à sua relação com triângulos retângulos e a aplicação direta do \textbf{Teorema de Pitágoras} durante a dedução trigonométrica destas derivadas.
  \end{itemize}
\end{notebox}


% ────────────────────────────────────────────────────────
\subsection{Exemplos Resolvidos}

\begin{exbox}{1 --- Derivada com Argumento Linear}
  Calcular a derivada da função $f(x) = \arcsin(3x)$.
\end{exbox}
\begin{solbox}
  Neste caso, o nosso argumento é $u(x) = 3x$. A sua derivada é $u'(x) = 3$.
  Aplicando a regra da cadeia para o arco-seno:
  \[
    \frac{\dd}{\dd x}\arcsin(3x) = \frac{(3x)'}{\sqrt{1-(3x)^2}} = \frac{3}{\sqrt{1-9x^2}}
  \]
  \textbf{Solução Final:} $\boxed{\frac{3}{\sqrt{1-9x^2}}}$
\end{solbox}

\begin{exbox}{2 --- Derivada com Argumento Quadrático}
  Calcular a derivada da função $f(x) = \arctan(x^2)$.
\end{exbox}
\begin{solbox}
  O argumento é $u(x) = x^2$, cuja derivada é $u'(x) = 2x$.
  Aplicando a regra da cadeia para o arco-tangente:
  \[
    \frac{\dd}{\dd x}\arctan(x^2) = \frac{(x^2)'}{1+(x^2)^2} = \frac{2x}{1+x^4}
  \]
  \textbf{Solução Final:} $\boxed{\frac{2x}{1+x^4}}$
\end{solbox}


% ────────────────────────────────────────────────────────
\subsection{Exercícios Práticos}

\begin{exercisebox}{1 --- Aplicação Direta da Regra da Cadeia}
  Calcule $\displaystyle \frac{\dd}{\dd x}\arccos(2x)$ usando a regra da cadeia.
\end{exercisebox}
\begin{solbox}
  Seja $u(x) = 2x$. A derivada de $u(x)$ é $2$.
  A fórmula da derivada do arco-cosseno é análoga à do arco-seno, mas com sinal negativo:
  \[
    \frac{\dd}{\dd x}\arccos(2x) = -\frac{(2x)'}{\sqrt{1-(2x)^2}} = -\frac{2}{\sqrt{1-4x^2}}
  \]
  \textbf{Solução Final:} $\boxed{-\frac{2}{\sqrt{1-4x^2}}}$
\end{solbox}

\begin{exercisebox}{2 --- Arco-tangente com Argumento Racional}
  Determine $\displaystyle \frac{\dd}{\dd x}\arctan\left(\frac{1}{x}\right)$.
\end{exercisebox}
\begin{solbox}
  Temos o argumento $u(x) = \frac{1}{x} = x^{-1}$. A sua derivada é $u'(x) = -x^{-2} = -\frac{1}{x^2}$.
  Substituindo na fórmula do $\arctan u$:
  \[
    \frac{\dd}{\dd x}\arctan\left(\frac{1}{x}\right) = \frac{-\frac{1}{x^2}}{1+\left(\frac{1}{x}\right)^2}
  \]
  Para simplificar a expressão algébrica, multiplicamos o numerador e o denominador por $x^2$:
  \[
    = \frac{-\frac{1}{x^2} \cdot x^2}{\left(1+\frac{1}{x^2}\right) \cdot x^2} = \frac{-1}{x^2+1} = -\frac{1}{1+x^2}
  \]
  \textbf{Nota curiosa:} O resultado é o oposto da derivada de $\arctan(x)$, o que se deve à identidade trigonométrica $\arctan(1/x) = \frac{\pi}{2} - \arctan(x)$ para $x>0$.
\end{solbox}

\begin{exercisebox}{3 --- Arco-seno com Argumento Cúbico}
  Calcule $\displaystyle \frac{\dd}{\dd x}\arcsin(x^3)$.
\end{exercisebox}
\begin{solbox}
  Identificamos o argumento como $u(x) = x^3$. A sua derivada é $u'(x) = 3x^2$.
  Aplicando a regra da cadeia para a derivada do $\arcsin$:
  \[
    \frac{\dd}{\dd x}\arcsin(x^3) = \frac{(x^3)'}{\sqrt{1-(x^3)^2}} = \frac{3x^2}{\sqrt{1-x^6}}
  \]
  \textbf{Solução Final:} $\boxed{\frac{3x^2}{\sqrt{1-x^6}}}$
\end{solbox}


% ────────────────────────────────────────────────────────
\subsection{Questões Propostas para Defesa Oral}

\begin{oralbox}
  \textbf{Q1.} Explique geometricamente porque é que a derivada de $\arcsin x$ envolve a existência de uma raiz quadrada no denominador. Que relação tem isto com o Teorema de Pitágoras?
\end{oralbox}

\begin{oralbox}
  \textbf{Q2.} O que distingue analiticamente o $\arcsin(x)$ do $\arccos(x)$ em termos da sua derivada? E geometricamente, como se comportam as retas tangentes a estas duas curvas?
\end{oralbox}

\begin{oralbox}
  \textbf{Q3.} Explique, por palavras suas, o papel fundamental da Regra da Cadeia na derivação de funções trigonométricas inversas compostas. O que falharia se derivássemos $\arctan(x^2)$ apenas como $\frac{1}{1+x^4}$?
\end{oralbox}

\begin{oralbox}
  \textbf{Q4.} A função $\arctan x$ é derivável em todos os números reais, enquanto o $\arcsin x$ e $\arccos x$ apenas são deriváveis no intervalo aberto $]-1, 1[$. Justifique matematicamente esta diferença observando os denominadores das suas derivadas.
\end{oralbox}
% ════════════════════════════════════════════════════════════════
% 1.3 - DERIVADAS DE ORDEM SUPERIOR E POLINÓMIO DE TAYLOR
% ════════════════════════════════════════════════════════════════

\section{Derivadas de Ordem Superior e Polinómio de Taylor}

% ────────────────────────────────────────────────────────
\subsection{Enquadramento Teórico}

\begin{notebox}[title={Intuição Geral}]
  Derivar uma função várias vezes sucessivas revela padrões profundos e comportamentos finos da curva geométrica:
  \begin{itemize}
      \item \textbf{1ª Derivada ($f'$):} Dá-nos a \emph{velocidade} ou taxa de variação (inclinação da reta tangente).
      \item \textbf{2ª Derivada ($f''$):} Dá-nos a \emph{aceleração} ou a taxa de variação da velocidade. Geometricamente, mede a \emph{curvatura} ou concavidade da função.
      \item \textbf{3ª Derivada em diante ($f'''$, $f^{(4)}$, \dots):} Refletem comportamentos de ordem superior (como o \emph{jerk} na física) e ajustes geométricos cada vez mais finos.
  \end{itemize}
  Juntando todas estas derivadas num único ponto, conseguimos construir o \textbf{Polinómio de Taylor}, que aproxima uma função complexa através de uma expressão polinomial simples.
\end{notebox}

\begin{defbox}[title={Derivadas Sucessivas}]
  Se $f$ é uma função derivável e a sua derivada $f'$ também o é, podemos continuar a derivar sucessivamente. Definimos rigorosamente:
  \[
    f'' = (f')', \qquad f''' = (f'')', \qquad f^{(n)} = \left(f^{(n-1)}\right)'
  \]
  A notação $f^{(n)}$ (com o $n$ entre parêntesis) é usada para representar a derivada de ordem $n$, evitando a confusão com potências algébricas ($f^n$).
\end{defbox}

\begin{defbox}[title={Polinómio de Taylor}]
  Seja $f$ uma função que admite derivadas até à ordem $n$ no ponto $x=a$. O \textbf{polinómio de Taylor de ordem $n$} de $f$ centrado em $a$ é dado por:
  \[
    T_n(x) = \sum_{k=0}^n \frac{f^{(k)}(a)}{k!}(x-a)^k
  \]
  Expandindo a notação de somatório, obtemos:
  \[
    T_n(x) = f(a) + f'(a)(x-a) + \frac{f''(a)}{2!}(x-a)^2 + \frac{f'''(a)}{3!}(x-a)^3 + \dots + \frac{f^{(n)}(a)}{n!}(x-a)^n
  \]
  \textit{Nota:} Quando o desenvolvimento é feito em torno da origem ($a=0$), o polinómio é frequentemente designado por \textbf{Polinómio de Maclaurin}.
\end{defbox}

% ────────────────────────────────────────────────────────
\subsection{Exemplos Resolvidos}

\begin{exbox}{1 --- Derivadas Sucessivas e Taylor da Exponencial}
  Derive a função $f(x)=e^x$ até à 4ª ordem e construa o respetivo polinómio de Taylor centrado em $a=0$.
\end{exbox}
\begin{solbox}
  A função exponencial de base $e$ é única porque é a sua própria derivada. Assim, todas as derivadas sucessivas são idênticas:
  \[
    f(x) = f'(x) = f''(x) = f'''(x) = f^{(4)}(x) = e^x
  \]
  Avaliando todas estas derivadas no ponto $a=0$, obtemos:
  \[
    f(0) = f'(0) = f''(0) = f'''(0) = f^{(4)}(0) = e^0 = 1
  \]
  Aplicando a fórmula do Polinómio de Taylor para $n=4$:
  \[
    T_4(x) = 1 + 1\cdot(x-0) + \frac{1}{2!}(x-0)^2 + \frac{1}{3!}(x-0)^3 + \frac{1}{4!}(x-0)^4
  \]
  \textbf{Solução Final:} $\boxed{T_4(x) = 1 + x + \frac{x^2}{2} + \frac{x^3}{6} + \frac{x^4}{24}}$
\end{solbox}

\begin{exbox}{2 --- Aproximação de Funções com Radicais}
  Aproximar a função $f(x) = \sqrt{1+x}$ através de um Polinómio de Taylor de ordem 3 perto da origem ($a=0$).
\end{exbox}
\begin{solbox}
  Reescrevemos a função como $f(x) = (1+x)^{1/2}$ e calculamos as derivadas com a regra da potência e da cadeia:
  \begin{align*}
      f(x)   &= (1+x)^{1/2} & \implies f(0) &= 1 \\
      f'(x)  &= \frac{1}{2}(1+x)^{-1/2} & \implies f'(0) &= \frac{1}{2} \\
      f''(x) &= -\frac{1}{4}(1+x)^{-3/2} & \implies f''(0) &= -\frac{1}{4} \\
      f'''(x)&= \frac{3}{8}(1+x)^{-5/2} & \implies f'''(0) &= \frac{3}{8}
  \end{align*}
  \begin{warnbox}[title={Atenção aos Fatoriais!}]
    Um erro muito comum é esquecer de dividir as derivadas pelos fatoriais da fórmula de Taylor.
  \end{warnbox}
  Aplicando a fórmula correta $T_3(x) = f(0) + f'(0)x + \frac{f''(0)}{2!}x^2 + \frac{f'''(0)}{3!}x^3$:
  \[
    T_3(x) = 1 + \frac{1}{2}x + \frac{-1/4}{2}x^2 + \frac{3/8}{6}x^3
  \]
  \textbf{Solução Final:} $\boxed{T_3(x) = 1 + \frac{1}{2}x - \frac{1}{8}x^2 + \frac{1}{16}x^3}$
\end{solbox}

% ────────────────────────────────────────────────────────
\subsection{Exercícios Práticos}

\begin{exercisebox}{1 --- Derivadas de Ordem Superior}
  Encontre as três primeiras derivadas sucessivas da função $f(x)=\ln(1+x)$.
\end{exercisebox}
\begin{solbox}
  Usando as regras de derivação (sabendo que a derivada de $\ln(u)$ é $u'/u$):
  \begin{align*}
      f'(x) &= \frac{1}{1+x} = (1+x)^{-1} \\[1ex]
      f''(x) &= -1 \cdot (1+x)^{-2} \cdot (1+x)' = -(1+x)^{-2} = -\frac{1}{(1+x)^2} \\[1ex]
      f'''(x) &= -(-2) \cdot (1+x)^{-3} \cdot 1 = 2(1+x)^{-3} = \frac{2}{(1+x)^3}
  \end{align*}
\end{solbox}

\begin{exercisebox}{2 --- Polinómio de Taylor Fora da Origem}
  Construa o polinómio de Taylor de ordem 2 da função $f(x) = \sin x$ centrado em $a=\pi/4$.
\end{exercisebox}
\begin{solbox}
  Calculamos a função e as suas duas primeiras derivadas no ponto $a = \pi/4$:
  \begin{align*}
      f(x) &= \sin x & \implies f(\pi/4) &= \frac{\sqrt{2}}{2} \\
      f'(x) &= \cos x & \implies f'(\pi/4) &= \frac{\sqrt{2}}{2} \\
      f''(x) &= -\sin x & \implies f''(\pi/4) &= -\frac{\sqrt{2}}{2}
  \end{align*}
  A fórmula de Taylor para $n=2$ é $T_2(x) = f(a) + f'(a)(x-a) + \frac{f''(a)}{2!}(x-a)^2$:
  \[
    T_2(x) = \frac{\sqrt{2}}{2} + \frac{\sqrt{2}}{2}\left(x-\frac{\pi}{4}\right) + \frac{-\sqrt{2}/2}{2}\left(x-\frac{\pi}{4}\right)^2
  \]
  \textbf{Solução Final:} $\boxed{T_2(x) = \frac{\sqrt{2}}{2} + \frac{\sqrt{2}}{2}\left(x-\frac{\pi}{4}\right) - \frac{\sqrt{2}}{4}\left(x-\frac{\pi}{4}\right)^2}$
\end{solbox}

\begin{exercisebox}{3 --- Série Geométrica via Taylor}
  Determine o Polinómio de Taylor $T_3(x)$ centrado em $a=0$ para a função $f(x)=\frac{1}{1-x}$.
\end{exercisebox}
\begin{solbox}
  Reescrevemos $f(x) = (1-x)^{-1}$ e derivamos usando a regra da cadeia (com $u' = -1$):
  \begin{align*}
      f(x) &= (1-x)^{-1} & \implies f(0) &= 1 \\
      f'(x) &= -(1-x)^{-2}(-1) = (1-x)^{-2} & \implies f'(0) &= 1 \\
      f''(x) &= -2(1-x)^{-3}(-1) = 2(1-x)^{-3} & \implies f''(0) &= 2 \\
      f'''(x) &= -6(1-x)^{-4}(-1) = 6(1-x)^{-4} & \implies f'''(0) &= 6
  \end{align*}
  Construindo o polinómio:
  \[
    T_3(x) = 1 + 1\cdot x + \frac{2}{2!}x^2 + \frac{6}{3!}x^3
  \]
  Simplificando os fatoriais ($2!=2$ e $3!=6$):
  \textbf{Solução Final:} $\boxed{T_3(x) = 1 + x + x^2 + x^3}$
\end{solbox}

% ────────────────────────────────────────────────────────
\subsection{Questões Propostas para Defesa Oral}

\begin{oralbox}
  \textbf{Q1.} Explique, intuitivamente, o que significa a segunda derivada ser positiva ou negativa num determinado ponto de uma curva geométrica. Qual é a relação disto com o comportamento das retas tangentes?
\end{oralbox}

\begin{oralbox}
  \textbf{Q2.} Em que situações prático-científicas ou computacionais é que a utilização do Polinómio de Taylor é extremamente útil face ao cálculo direto da função original? Dê um exemplo.
\end{oralbox}

\begin{oralbox}
  \textbf{Q3.} Por que motivo é que a função exponencial $e^x$ é muitas vezes considerada a função "mais simples" para derivar repetidamente e para construir as séries de Taylor?
\end{oralbox}
% ════════════════════════════════════════════════════════════════
% 1.4 - CURVAS PARAMÉTRICAS E DERIVAÇÃO
% ════════════════════════════════════════════════════════════════

\section{Curvas Paramétricas e Derivação}

% ────────────────────────────────────────────────────────
\subsection{Enquadramento Teórico}

\begin{notebox}[title={Intuição Geral}]
  As curvas paramétricas são a ferramenta ideal para descrever \textbf{movimentos} contínuos. Em vez de forçarmos uma variável a depender da outra (como em $y = f(x)$), expressamos ambas as coordenadas como funções de uma terceira variável independente: o parâmetro $t$, que é frequentemente interpretado como o ``tempo''.
  Assim, a posição de um objeto a cada instante é dada pelo ponto $(x(t), y(t))$. A \emph{derivada} destas funções indica-nos exatamente como a posição está a mudar, ou seja, dá-nos a \textbf{velocidade} do objeto nesse instante.
\end{notebox}

\begin{defbox}[title={Curva Paramétrica e Vetor Velocidade}]
  Uma \textbf{curva paramétrica} no plano bidimensional é definida por um par de funções contínuas:
  \[
    \gamma(t) = (x(t), y(t)), \quad t \in I
  \]
  onde $I$ é um intervalo de números reais.
  
  Se as funções $x(t)$ e $y(t)$ forem deriváveis, a \textbf{derivada da curva} (ou vetor velocidade) é obtida derivando cada componente individualmente em ordem a $t$:
  \[
    \gamma'(t) = (x'(t), y'(t))
  \]
  Este vetor $\gamma'(t)$ é geometricamente \textbf{tangente} à curva geométrica no ponto $\gamma(t)$.
\end{defbox}

\begin{defbox}[title={Comprimento de uma Curva (Arco)}]
  Se o vetor velocidade for contínuo e a curva não se cruzar a si própria, o \textbf{comprimento total} $L$ do percurso percorrido entre o instante $t=a$ e $t=b$ é dado pelo integral da norma (tamanho) do vetor velocidade:
  \[
    L = \int_a^b \|\gamma'(t)\| \,dt = \int_a^b \sqrt{(x'(t))^2 + (y'(t))^2} \,dt
  \]
\end{defbox}

% ────────────────────────────────────────────────────────
\subsection{Exemplos Resolvidos}

\begin{exbox}{1 --- A Circunferência Padrão}
  Considere a curva paramétrica geométrica que descreve uma circunferência de raio 1:
  \[
    \gamma(t) = (\cos t, \sin t)
  \]
  Calcule a sua derivada e interprete o resultado geometricamente.
\end{exbox}
\begin{solbox}
  Derivando componente a componente:
  \[
    \gamma'(t) = (-\sin t, \cos t)
  \]
  \textbf{Interpretação:} O vetor velocidade $\gamma'(t)$ é sempre ortogonal ao vetor posição $\gamma(t)$ (o produto interno entre eles é $-\sin t \cos t + \cos t \sin t = 0$), o que confirma que a velocidade é sempre \textbf{tangente} ao círculo. Além disso, a norma do vetor velocidade é:
  \[
    \|\gamma'(t)\| = \sqrt{(-\sin t)^2 + (\cos t)^2} = \sqrt{\sin^2 t + \cos^2 t} = \sqrt{1} = 1
  \]
  Como a norma é constante e igual a 1, trata-se de um \textbf{movimento circular uniforme}.
\end{solbox}

\begin{exbox}{2 --- Comprimento de uma Espiral Simples}
  Determine a expressão integral para o comprimento da espiral definida por $\gamma(t)=(t\cos t, t\sin t)$ entre $t=0$ e $t=a$.
\end{exbox}
\begin{solbox}
  Primeiro, determinamos o vetor velocidade usando a regra do produto em cada componente:
  \begin{align*}
      x'(t) &= (\cos t - t\sin t) \\
      y'(t) &= (\sin t + t\cos t)
  \end{align*}
  O comprimento entre $0$ e $a$ é dado por:
  \[
    L = \int_0^a \sqrt{(\cos t - t\sin t)^2 + (\sin t + t\cos t)^2} \,dt
  \]
  \textbf{Simplificação do interior da raiz:}
  Expandindo os quadrados perfeitos:
  \begin{align*}
      (x')^2 &= \cos^2 t - 2t\cos t\sin t + t^2\sin^2 t \\
      (y')^2 &= \sin^2 t + 2t\sin t\cos t + t^2\cos^2 t
  \end{align*}
  Somando as duas expressões, o termo cruzado anula-se:
  \[
    (x')^2 + (y')^2 = (\cos^2 t + \sin^2 t) + t^2(\sin^2 t + \cos^2 t) = 1 + t^2
  \]
  \textbf{Resultado final:}
  A expressão integral simplificada para o comprimento da espiral é:
  \[
    \boxed{ L = \int_0^a \sqrt{1+t^2} \,dt }
  \]
\end{solbox}

% ────────────────────────────────────────────────────────
\subsection{Exercícios Práticos}

\begin{exercisebox}{1 --- Derivação Paramétrica Simples}
  Derive a curva polinomial $\gamma(t)=(t^2, t^3)$.
\end{exercisebox}
\begin{solbox}
  Basta derivar cada uma das coordenadas em ordem à variável independente $t$:
  \begin{itemize}
      \item $x(t) = t^2 \implies x'(t) = 2t$
      \item $y(t) = t^3 \implies y'(t) = 3t^2$
  \end{itemize}
  \textbf{Solução Final:} O vetor velocidade é $\boxed{\gamma'(t) = (2t, 3t^2)}$.
\end{solbox}

\begin{exercisebox}{2 --- Cálculo Efetivo do Comprimento de Arco}
  Calcule o comprimento da curva linear $\gamma(t)=(3t, 4t)$ entre $t=0$ e $t=1$.
\end{exercisebox}
\begin{solbox}
  Primeiro, calculamos o vetor velocidade:
  \[
    \gamma'(t) = (3, 4)
  \]
  A norma deste vetor (a velocidade escalar) é:
  \[
    \|\gamma'(t)\| = \sqrt{3^2 + 4^2} = \sqrt{9 + 16} = \sqrt{25} = 5
  \]
  O comprimento total é o integral desta norma entre os instantes fornecidos:
  \[
    L = \int_0^1 5 \,dt = \Big[ 5t \Big]_0^1 = 5(1) - 5(0) = \mathbf{5}
  \]
  \textit{Nota geométrica:} Esta curva descreve um segmento de reta que vai do ponto $(0,0)$ ao ponto $(3,4)$. A distância direta é precisamente 5 (conforme o Teorema de Pitágoras).
\end{solbox}

\begin{exercisebox}{3 --- Ortogonalidade do Movimento}
  Mostre que a curva $\gamma(t)=(e^t, e^{-t})$ tem um vetor velocidade perpendicular ao seu \emph{vetor posição} no instante $t=0$.
\end{exercisebox}
\begin{solbox}
  Para provar a perpendicularidade (ortogonalidade), o produto interno (produto escalar) entre o vetor posição e o vetor velocidade em $t=0$ tem de ser igual a zero.
  
  \textbf{1. Vetor Posição em $t=0$:}
  \[
    \gamma(0) = (e^0, e^{-0}) = (1, 1)
  \]
  
  \textbf{2. Vetor Velocidade em $t=0$:}
  Derivamos a curva original:
  \[
    \gamma'(t) = (e^t, -e^{-t})
  \]
  Avaliando no instante $t=0$:
  \[
    \gamma'(0) = (e^0, -e^{-0}) = (1, -1)
  \]
  
  \textbf{3. Produto Interno:}
  Fazendo o produto escalar entre $\gamma(0)$ e $\gamma'(0)$:
  \[
    \gamma(0) \cdot \gamma'(0) = (1)(1) + (1)(-1) = 1 - 1 = \mathbf{0}
  \]
  Como o produto escalar é zero, os vetores formam um ângulo de $90^\circ$. A afirmação está provada. \checkmark
\end{solbox}

% ────────────────────────────────────────────────────────
\subsection{Questões Propostas para Defesa Oral}

\begin{oralbox}
  \textbf{Q1.} Se tivesse de explicar o conceito de "curva paramétrica" a alguém sem formação matemática avançada, que analogia do quotidiano usaria? Como explicaria o papel do parâmetro $t$?
\end{oralbox}

\begin{oralbox}
  \textbf{Q2.} Num movimento parametrizado, o que representa geometricamente e fisicamente o \emph{vetor velocidade}? Como é que ele se relaciona com a curva principal?
\end{oralbox}

\begin{oralbox}
  \textbf{Q3.} Justifique a fórmula do comprimento de uma curva. Por que razão esta fórmula envolve especificamente a raiz quadrada da soma dos quadrados das derivadas? A que teorema célebre isto remete?
\end{oralbox}

\begin{oralbox}
  \textbf{Q4.} Qual é a grande vantagem (ou diferença estrutural) de parametrizar uma curva com $t$, em vez de desenhá-la estaticamente através de uma função comum $y = f(x)$? Dê um exemplo de uma forma geométrica que a parametrização resolve facilmente, mas que $y=f(x)$ não consegue.
\end{oralbox}

% ---------------------------------------------------------
% CAPÍTULO 2
% ---------------------------------------------------------
% ════════════════════════════════════════════════════════════════
% CAPÍTULO 2 -- CÁLCULO INTEGRAL EM R
% Ficheiro para \include{} no template principal.
% ════════════════════════════════════════════════════════════════

\chapter{Cálculo Integral em \texorpdfstring{$\mathbb{R}$}{R}}

Este capítulo segue os tópicos do programa 2.1--2.5 especificados para os
cursos de EI/EE e EM. Abrange primitivas e as suas técnicas de cálculo,
o integral definido, o Teorema Fundamental do Cálculo Integral e
aplicações ao cálculo de áreas, volumes e comprimentos de arco.

% ────────────────────────────────────────────────────────
\section{Primitivas e Técnicas de Primitivação}

% ────────────────────────────────────────────────────────
\subsection{Enquadramento Teórico}

\begin{defbox}[title={Primitiva (Antiderivada)}]
  Uma função $F : I \to \mathbb{R}$ diz-se uma \emph{primitiva} (ou
  \emph{antiderivada}) de $f : I \to \mathbb{R}$ num intervalo aberto $I$ se
  \[
    F'(x) = f(x) \quad \forall\, x \in I.
  \]
  A família de todas as primitivas escreve-se $\displaystyle\int f(x)\,dx = F(x)+C$,
  onde $C \in \mathbb{R}$ é uma constante arbitrária.
\end{defbox}

\begin{notebox}
  Duas primitivas da mesma função diferem por uma constante, pelo que o
  integral indefinido é único a menos de $C$.
\end{notebox}

\subsubsection*{Tabela de Primitivas Imediatas}

\begin{center}
\renewcommand{\arraystretch}{1.5}
\begin{tabular}{cc}
  \toprule
  $f(x)$ & $\displaystyle\int f(x)\,dx$ \\
  \midrule
  $x^n\ (n\neq-1)$ & $\dfrac{x^{n+1}}{n+1}+C$ \\[6pt]
  $\dfrac{1}{x}$ & $\ln|x|+C$ \\[6pt]
  $e^x$ & $e^x+C$ \\[6pt]
  $\sin x$ & $-\cos x+C$ \\[6pt]
  $\cos x$ & $\sin x+C$ \\[6pt]
  $\sec^2 x$ & $\tan x+C$ \\[6pt]
  $\dfrac{1}{1+x^2}$ & $\arctan x + C$ \\[6pt]
  $\dfrac{1}{\sqrt{1-x^2}}$ & $\arcsin x + C$ \\
  \bottomrule
\end{tabular}
\end{center}

\subsubsection*{Primitivação por Partes}

\begin{thmbox}[title={Fórmula de Primitivação por Partes}]
\[
  \int u\,dv = uv - \int v\,du.
\]
Escolhe-se $u$ e $dv$ de modo que $\int v\,du$ seja mais simples. O
mnemónico \textbf{LIATE} (Logarítmica, Inversa trigonométrica, Algébrica,
Trigonométrica, Exponencial) sugere a prioridade na escolha de $u$.
\end{thmbox}

\subsubsection*{Primitivação por Substituição}

\begin{thmbox}[title={Fórmula de Substituição}]
Se $x = g(t)$, então
\[
  \int f(x)\,dx = \int f(g(t))\,g'(t)\,dt.
\]
\end{thmbox}

\subsubsection*{Potências de Funções Trigonométricas}

Para integrar $\sin^m x\cos^n x$:
\begin{itemize}
  \item Se $m$ for ímpar: substitua $u=\cos x$.
  \item Se $n$ for ímpar: substitua $u=\sin x$.
  \item Se ambos forem pares: use as identidades do ângulo duplo
        $\sin^2 x = \tfrac{1-\cos 2x}{2}$,
        $\cos^2 x = \tfrac{1+\cos 2x}{2}$.
\end{itemize}

\subsubsection*{Frações Racionais (Decomposição em Frações Parciais)}

Toda a função racional própria $\dfrac{P(x)}{Q(x)}$ pode ser decomposta em
frações parciais. Para uma raíz real simples $r$ de $Q$:
\[
  \frac{A}{x-r}.
\]
Para um fator quadrático irredutível $x^2+px+q$:
\[
  \frac{Bx+C}{x^2+px+q}.
\]

% ────────────────────────────────────────────────────────
\subsection{Exemplos Resolvidos}

\begin{exbox}{1 --- Primitivação por Partes}
  Calcule $\displaystyle\int x e^x\,dx$.
\end{exbox}
\begin{solbox}
  Tome-se $u=x$, $dv=e^x\,dx$. Então $du=dx$, $v=e^x$.
  \[
    \int x e^x\,dx = xe^x - \int e^x\,dx = xe^x - e^x + C = e^x(x-1)+C.
  \]
\end{solbox}

\begin{exbox}{2 --- Substituição}
  Calcule $\displaystyle\int \frac{2x}{x^2+1}\,dx$.
\end{exbox}
\begin{solbox}
  Seja $u = x^2+1$, logo $du = 2x\,dx$.
  \[
    \int \frac{2x}{x^2+1}\,dx = \int \frac{du}{u} = \ln|u|+C = \ln(x^2+1)+C.
  \]
\end{solbox}

\begin{exbox}{3 --- Potências de Funções Trigonométricas}
  Calcule $\displaystyle\int \sin^3 x\,dx$.
\end{exbox}
\begin{solbox}
  Escreva $\sin^3 x = \sin x(1-\cos^2 x)$; seja $u=\cos x$:
  \[
    \int \sin^3 x\,dx = -\int (1-u^2)\,du = -u + \tfrac{u^3}{3}+C
      = -\cos x + \tfrac{\cos^3 x}{3}+C.
  \]
\end{solbox}

\begin{exbox}{4 --- Frações Racionais}
  Calcule $\displaystyle\int \frac{1}{x^2-1}\,dx$.
\end{exbox}
\begin{solbox}
  Decomposição: $\dfrac{1}{x^2-1}=\dfrac{1}{2}\!\left(\dfrac{1}{x-1}-\dfrac{1}{x+1}\right)$.
  \[
    \int \frac{dx}{x^2-1} = \frac{1}{2}\ln\left|\frac{x-1}{x+1}\right|+C.
  \]
\end{solbox}


% ────────────────────────────────────────────────────────
\section{O Integral Definido}

% ────────────────────────────────────────────────────────
\subsection{Enquadramento Teórico}

\begin{defbox}[title={Integral de Riemann}]
  Seja $f:[a,b]\to\mathbb{R}$ limitada. Para uma partição
  $P=\{a=x_0<x_1<\cdots<x_n=b\}$, a \emph{soma de Riemann} é
  \[
    S(f,P) = \sum_{i=1}^n f(c_i)(x_i - x_{i-1}),\quad c_i\in[x_{i-1},x_i].
  \]
  Se $S(f,P)$ convergir para um limite $L$ quando a norma da partição
  $\|P\|\to 0$ (independentemente da escolha de $c_i$), então $f$ é
  \emph{integrável no sentido de Riemann} e $\displaystyle\int_a^b f(x)\,dx = L$.
\end{defbox}

\begin{thmbox}[title={Propriedades Principais do Integral Definido}]
\begin{align*}
  &\int_a^b [\alpha f+\beta g]\,dx = \alpha\int_a^b f\,dx + \beta\int_a^b g\,dx
    \quad (\text{linearidade})\\[4pt]
  &\int_a^b f\,dx = \int_a^c f\,dx + \int_c^b f\,dx
    \quad (\text{aditividade em } [a,c]\cup[c,b])\\[4pt]
  &f(x)\geq 0 \implies \int_a^b f\,dx \geq 0
    \quad (\text{positividade})
\end{align*}
\end{thmbox}


% ────────────────────────────────────────────────────────
\section{O Teorema Fundamental do Cálculo Integral}

% ────────────────────────────────────────────────────────
\subsection{Enquadramento Teórico}

\begin{thmbox}[title={Teorema Fundamental do Cálculo Integral}]
  Seja $f:[a,b]\to\mathbb{R}$ contínua.
  \begin{enumerate}
    \item \textbf{(Parte 1)} A função $G(x)=\displaystyle\int_a^x f(t)\,\dd t$
          é diferenciável em $(a,b)$ e $G'(x)=f(x)$.
    \item \textbf{(Parte 2)} Se $F$ for qualquer primitiva de $f$, então
          \[
            \int_a^b f(x)\,\dd x = F(b)-F(a) =: \Big[F(x)\Big]_a^b.
          \]
  \end{enumerate}
\end{thmbox}

% ────────────────────────────────────────────────────────
\subsection{Exemplos Resolvidos}

\begin{exbox}{5 --- Aplicação do Teorema Fundamental}
  Calcule $\displaystyle\int_0^{\pi} \sin x\,\dd x$.
\end{exbox}
\begin{solbox}
  Uma primitiva de $\sin x$ é $-\cos x$.
  \[
    \int_0^{\pi}\sin x\,\dd x = \Big[-\cos x\Big]_0^{\pi}
      = (-\cos\pi)-(-\cos 0) = 1+1 = 2.
  \]
\end{solbox}


% ────────────────────────────────────────────────────────
\section{Aplicações e Comprimento do Arco}

% ────────────────────────────────────────────────────────
\subsection{Enquadramento Teórico}

\subsubsection*{Área entre Curvas}

Se $f(x)\geq g(x)$ em $[a,b]$, a área da região delimitada é:
\[
  A = \int_a^b \bigl[f(x)-g(x)\bigr]\,\dd x.
\]

\subsubsection*{Volume de um Sólido de Revolução (Método dos Discos)}

Ao rodar $y=f(x)\geq 0$ em torno do eixo $Ox$ de $a$ a $b$:
\[
  V = \pi\int_a^b [f(x)]^2\,\dd x.
\]

\subsubsection*{Comprimento do Arco}

O comprimento da curva $y=f(x)$ em $[a,b]$:
\[
  L = \int_a^b \sqrt{1+[f'(x)]^2}\,\dd x.
\]

Para uma curva paramétrica $(x(t),y(t))$, $t\in[\alpha,\beta]$:
\[
  L = \int_\alpha^\beta \sqrt{[\dot x(t)]^2+[\dot y(t)]^2}\,\dd t.
\]

% ────────────────────────────────────────────────────────
\subsection{Exemplos Resolvidos}

\begin{exbox}{6 --- Área entre Curvas}
  Determine a área entre $f(x)=x^2$ e $g(x)=x$ em $[0,1]$.
\end{exbox}
\begin{solbox}
  Como $x\geq x^2$ em $[0,1]$:
  \[
    A = \int_0^1(x-x^2)\,\dd x = \left[\frac{x^2}{2}-\frac{x^3}{3}\right]_0^1
      = \frac{1}{2}-\frac{1}{3} = \frac{1}{6}.
  \]
\end{solbox}

\begin{exbox}{7 --- Volume de uma Esfera}
  Volume da esfera de raio $r$: roda-se $f(x)=\sqrt{r^2-x^2}$ em
  $[-r,r]$ em torno do eixo $Ox$.
\end{exbox}
\begin{solbox}
  \textbf{Passo 1 -- Identificar a função.} A semicircunferência superior de
  raio $r$ é dada por $f(x)=\sqrt{r^2-x^2}$, com $x\in[-r,r]$. Ao rodar
  este arco em torno do eixo $Ox$ obtém-se uma esfera completa.

  \textbf{Passo 2 -- Aplicar a fórmula dos discos.}
  \[
    V = \pi\int_{-r}^r [f(x)]^2\,\dd x = \pi\int_{-r}^r (r^2-x^2)\,\dd x.
  \]

  \textbf{Passo 3 -- Calcular o integral.}
  \[
    V = \pi\left[r^2 x - \frac{x^3}{3}\right]_{-r}^r
      = \pi\left[\left(r^3-\frac{r^3}{3}\right)-\left(-r^3+\frac{r^3}{3}\right)\right]
      = \pi\cdot\frac{4r^3}{3} = \frac{4}{3}\pi r^3.
  \]
\end{solbox}

\begin{exbox}{8 --- Comprimento do Arco}
  Determine o comprimento de $y=\dfrac{2}{3}x^{3/2}$ de $x=0$ a $x=3$.
\end{exbox}
\begin{solbox}
  $f'(x)=x^{1/2}$, logo $1+[f']^2=1+x$.
  \[
    L = \int_0^3\sqrt{1+x}\,\dd x = \left[\tfrac{2}{3}(1+x)^{3/2}\right]_0^3
      = \tfrac{2}{3}(8-1) = \tfrac{14}{3}.
  \]
\end{solbox}

% ────────────────────────────────────────────────────────
\subsection{Exercícios Práticos}

\begin{exercisebox}{1}
  Calcule os seguintes integrais indefinidos:
  \begin{enumerate}[label=(\alph*)]
    \item $\displaystyle\int (3x^2-2x+5)\,\dd x$
    \item $\displaystyle\int \frac{x}{x^2+4}\,\dd x$
    \item $\displaystyle\int x^2 e^x\,\dd x$
    \item $\displaystyle\int \frac{x+1}{x^2+x-2}\,\dd x$
    \item $\displaystyle\int \cos^2 x\,\dd x$
  \end{enumerate}
\end{exercisebox}
\begin{solbox}
  \textbf{(a)} Aplica-se a linearidade e a regra da potência:
  \[
    \int (3x^2-2x+5)\,\dd x = 3\cdot\frac{x^3}{3} - 2\cdot\frac{x^2}{2} + 5x + C
    = x^3 - x^2 + 5x + C.
  \]

  \textbf{(b)} Reconhece-se a forma $\displaystyle\int \frac{f'(x)}{f(x)}\,\dd x$.
  Seja $u=x^2+4$, $du=2x\,\dd x$, logo $x\,\dd x=\dfrac{du}{2}$:
  \[
    \int \frac{x}{x^2+4}\,\dd x = \frac{1}{2}\int\frac{du}{u}
    = \frac{1}{2}\ln|u|+C = \frac{1}{2}\ln(x^2+4)+C.
  \]

  \textbf{(c)} Aplica-se primitivação por partes duas vezes (LIATE: $u=x^2$,
  $dv=e^x\,\dd x$):
  \begin{align*}
    &\textbf{1.ª vez:}\quad u=x^2,\ dv=e^x\,\dd x \Rightarrow du=2x\,\dd x,\ v=e^x.\\
    &\int x^2 e^x\,\dd x = x^2 e^x - 2\int x e^x\,\dd x.\\
    &\textbf{2.ª vez:}\quad u=x,\ dv=e^x\,\dd x \Rightarrow du=\dd x,\ v=e^x.\\
    &\int x e^x\,\dd x = xe^x - e^x + C.\\
    &\textbf{Resultado:}\quad \int x^2 e^x\,\dd x = x^2 e^x - 2(xe^x - e^x)+C
      = e^x(x^2-2x+2)+C.
  \end{align*}

  \textbf{(d)} Fatoriza-se o denominador: $x^2+x-2=(x+2)(x-1)$.
  Decomposição em frações parciais:
  \[
    \frac{x+1}{(x+2)(x-1)} = \frac{A}{x+2}+\frac{B}{x-1}.
  \]
  Multiplicando ambos os membros por $(x+2)(x-1)$: $x+1 = A(x-1)+B(x+2)$.
  \begin{itemize}
    \item $x=1$: $2=3B \Rightarrow B=\tfrac{2}{3}$.
    \item $x=-2$: $-1=-3A \Rightarrow A=\tfrac{1}{3}$.
  \end{itemize}
  \[
    \int \frac{x+1}{x^2+x-2}\,\dd x = \frac{1}{3}\ln|x+2|+\frac{2}{3}\ln|x-1|+C.
  \]

  \textbf{(e)} Usa-se a identidade do ângulo duplo: $\cos^2 x = \dfrac{1+\cos 2x}{2}$.
  \[
    \int \cos^2 x\,\dd x = \int \frac{1+\cos 2x}{2}\,\dd x
    = \frac{x}{2} + \frac{\sin 2x}{4} + C.
  \]
\end{solbox}

\begin{exercisebox}{2}
  Calcule os seguintes integrais definidos:
  \begin{enumerate}[label=(\alph*)]
    \item $\displaystyle\int_1^e \ln x\,\dd x$
    \item $\displaystyle\int_0^1 x\sqrt{1-x^2}\,\dd x$
    \item $\displaystyle\int_0^{\pi/2}\sin^2 x\,\dd x$
  \end{enumerate}
\end{exercisebox}
\begin{solbox}
  \textbf{(a)} Primitivação por partes: $u=\ln x$, $dv=\dd x$, logo
  $du=\dfrac{1}{x}\,\dd x$, $v=x$.
  \[
    \int_1^e \ln x\,\dd x = \Big[x\ln x\Big]_1^e - \int_1^e x\cdot\frac{1}{x}\,\dd x
    = \Big[x\ln x\Big]_1^e - \Big[x\Big]_1^e
    = (e\cdot 1 - 1\cdot 0)-(e-1) = 1.
  \]

  \textbf{(b)} Substituição: $u=1-x^2$, $du=-2x\,\dd x$.
  Limites: $x=0\Rightarrow u=1$; $x=1\Rightarrow u=0$.
  \[
    \int_0^1 x\sqrt{1-x^2}\,\dd x = -\frac{1}{2}\int_1^0 \sqrt{u}\,\dd u
    = \frac{1}{2}\int_0^1 u^{1/2}\,\dd u
    = \frac{1}{2}\left[\frac{2}{3}u^{3/2}\right]_0^1 = \frac{1}{3}.
  \]

  \textbf{(c)} Identidade: $\sin^2 x = \dfrac{1-\cos 2x}{2}$.
  \[
    \int_0^{\pi/2}\sin^2 x\,\dd x = \int_0^{\pi/2}\frac{1-\cos 2x}{2}\,\dd x
    = \left[\frac{x}{2}-\frac{\sin 2x}{4}\right]_0^{\pi/2}
    = \frac{\pi}{4}.
  \]
\end{solbox}

\begin{exercisebox}{3}
  Determine a área da região delimitada por $y=x^3$ e $y=x$.
\end{exercisebox}
\begin{solbox}
  \textbf{Passo 1 -- Pontos de interseção.} $x^3=x \Rightarrow x(x^2-1)=0
  \Rightarrow x\in\{-1,0,1\}$.

  \textbf{Passo 2 -- Sinal da diferença.} Em $(-1,0)$: $x^3-x = x(x^2-1)$.
  Teste em $x=-\tfrac{1}{2}$: $(-\tfrac{1}{2})({\tfrac{1}{4}-1})>0$,
  logo $x^3\geq x$ em $(-1,0)$.
  Em $(0,1)$: teste em $x=\tfrac{1}{2}$: resultado negativo, logo $x\geq x^3$.

  \textbf{Passo 3 -- Cálculo por simetria.} A função $x^3-x$ é ímpar, pelo que
  a área em $(-1,0)$ iguala a área em $(0,1)$:
  \[
    A = 2\int_0^1(x-x^3)\,\dd x = 2\left[\frac{x^2}{2}-\frac{x^4}{4}\right]_0^1
    = 2\cdot\frac{1}{4} = \frac{1}{2}.
  \]
\end{solbox}

\begin{exercisebox}{4}
  Determine o volume do sólido obtido ao rodar a região delimitada por
  $y=\sqrt{x}$, $x=4$ e $y=0$ em torno do eixo $Ox$.
\end{exercisebox}
\begin{solbox}
  \textbf{Passo 1 -- Identificar o método.} A região é delimitada por baixo por
  $y=0$ e por cima por $y=\sqrt{x}$, com $x\in[0,4]$. Aplica-se o método dos
  discos.

  \textbf{Passo 2 -- Aplicar a fórmula.}
  \[
    V = \pi\int_0^4 [\sqrt{x}]^2\,\dd x = \pi\int_0^4 x\,\dd x.
  \]

  \textbf{Passo 3 -- Calcular.}
  \[
    V = \pi\left[\frac{x^2}{2}\right]_0^4 = \pi\cdot\frac{16}{2} = 8\pi.
  \]
\end{solbox}

\begin{exercisebox}{5}
  Calcule o comprimento de arco da curva $y = \ln(\cos x)$ de $x=0$ a
  $x=\pi/4$.
\end{exercisebox}
\begin{solbox}
  \textbf{Passo 1 -- Derivada.}
  \[
    y' = \frac{-\sin x}{\cos x} = -\tan x.
  \]

  \textbf{Passo 2 -- Expressão sob a raiz.}
  \[
    1+[y']^2 = 1+\tan^2 x = \sec^2 x.
  \]

  \textbf{Passo 3 -- Integral do comprimento.}
  \[
    L = \int_0^{\pi/4}\sqrt{\sec^2 x}\,\dd x = \int_0^{\pi/4}|\sec x|\,\dd x
    = \int_0^{\pi/4}\sec x\,\dd x.
  \]
  (Note-se que $\sec x > 0$ em $[0,\pi/4]$.)

  \textbf{Passo 4 -- Calcular $\displaystyle\int\sec x\,\dd x$.}
  Usando a primitiva conhecida $\displaystyle\int\sec x\,\dd x = \ln|\sec x+\tan x|+C$:
  \[
    L = \Big[\ln|\sec x + \tan x|\Big]_0^{\pi/4}
    = \ln\!\left(\sqrt{2}+1\right) - \ln(1+0)
    = \ln(\sqrt{2}+1).
  \]
\end{solbox}

% ────────────────────────────────────────────────────────
\subsection{Questões Propostas para Defesa Oral}

\begin{oralbox}
  \textbf{Q1.} Enuncie e demonstre o Teorema Fundamental do Cálculo Integral
  (Parte 2). Qual o papel da continuidade?
\end{oralbox}

\begin{oralbox}
  \textbf{Q2.} Explique a diferença entre integral indefinido e integral definido.
  Por que razão a constante arbitrária $C$ desaparece no caso definido?
\end{oralbox}

\begin{oralbox}
  \textbf{Q3.} Quando é que a primitivação por partes é mais útil? Dê um exemplo
  em que seja necessário aplicá-la duas vezes.
\end{oralbox}

\begin{oralbox}
  \textbf{Q4.} Descreva geometricamente o que representa o integral definido
  $\displaystyle\int_a^b f(x)\,\dd x$ quando $f$ toma valores negativos em parte de $[a,b]$.
\end{oralbox}

\begin{oralbox}
  \textbf{Q5.} Deduza a fórmula do comprimento de arco a partir de primeiros
  princípios, recorrendo ao Teorema de Pitágoras.
\end{oralbox}

\begin{oralbox}
  \textbf{Q6.} Compare o método dos discos e o método das cascas para o cálculo
  de volumes. Em que circunstâncias se prefere cada um?
\end{oralbox}

\begin{oralbox}
  \textbf{Q7.} Qual a relação entre a área sob uma curva e uma soma de Riemann?
  Descreva a convergência.
\end{oralbox}

% ---------------------------------------------------------
% CAPÍTULO 3
% ---------------------------------------------------------
% ════════════════════════════════════════════════════════════════
% CAPÍTULO 3 -- FUNÇÕES REAIS DE VÁRIAS VARIÁVEIS REAIS
% Ficheiro para \include{} no template principal.
% ════════════════════════════════════════════════════════════════

\chapter{Funções Reais de Várias Variáveis Reais}

Este capítulo abrange os tópicos do programa 3.1--3.5, incluindo domínio e
representação gráfica, limites e continuidade, derivadas parciais,
diferenciabilidade, regra da cadeia e extremos livres.

% ────────────────────────────────────────────────────────
\section{Definição, Domínio, Curvas de Nível e Representação Gráfica}

% ────────────────────────────────────────────────────────
\subsection{Enquadramento Teórico}

\begin{defbox}[title={Função Real de Várias Variáveis}]
  Uma \emph{função real de $n$ variáveis reais} é uma aplicação
  $f: D \subseteq \mathbb{R}^n \to \mathbb{R}$.
  Para $n=2$ escrevemos $f(x,y)$ e o seu \emph{gráfico} é a superfície
  \[
    \{(x,y,z)\in\mathbb{R}^3 \mid z = f(x,y),\, (x,y)\in D\}.
  \]
\end{defbox}

\begin{defbox}[title={Curva de Nível (Linha de Contorno)}]
  Para $f:D\subseteq\mathbb{R}^2\to\mathbb{R}$ e uma constante $c\in\mathbb{R}$,
  a \emph{curva de nível} (linha de contorno) à altura $c$ é
  \[
    L_c = \{(x,y)\in D \mid f(x,y) = c\}.
  \]
\end{defbox}

\subsubsection*{Representação Gráfica}

O gráfico de $f(x,y)$ é uma \emph{superfície} no espaço $\mathbb{R}^3$.
Para o representar, os métodos mais comuns são:

\begin{itemize}
  \item \textbf{Curvas de nível:} traçam-se as curvas $f(x,y)=c$ no plano
        $Oxy$ para vários valores de $c$, obtendo uma vista ``de cima'' da
        superfície. Curvas muito próximas indicam grande variação de $f$;
        curvas afastadas indicam variação suave.
  \item \textbf{Cortes verticais:} fixa-se $x=a$ ou $y=b$ e traça-se a
        curva resultante em $\mathbb{R}^2$, obtendo uma ``fatia'' da superfície.
\end{itemize}

% ────────────────────────────────────────────────────────
\subsection{Exemplos Resolvidos}

\begin{exbox}{1 --- Curvas de Nível de $f(x,y)=x^2+y^2$}
  Descreva as curvas de nível de $f(x,y)=x^2+y^2$.
\end{exbox}
\begin{solbox}
  A equação $f(x,y)=c$ escreve-se $x^2+y^2=c$.
  \begin{itemize}
    \item \textbf{Passo 1:} Se $c<0$, a equação $x^2+y^2=c$ não tem solução
          real (soma de quadrados nunca é negativa), pelo que a curva de nível
          é vazia.
    \item \textbf{Passo 2:} Se $c=0$, a única solução é $x=y=0$, ou seja,
          a curva de nível é o ponto $(0,0)$.
    \item \textbf{Passo 3:} Se $c>0$, a equação $x^2+y^2=c$ representa
          uma circunferência centrada na origem com raio $\sqrt{c}$.
  \end{itemize}
  Conclusão: as curvas de nível são circunferências concêntricas de raio
  $\sqrt{c}$ (para $c>0$), o ponto origem ($c=0$), ou vazias ($c<0$).
\end{solbox}

\begin{exbox}{2 --- Representação Gráfica}
  Descreva o gráfico de $f(x,y)=4-x^2-y^2$.
\end{exbox}
\begin{solbox}
  \textbf{Passo 1 -- Identificar o tipo de superfície.} A equação
  $z=4-x^2-y^2$ pode escrever-se $x^2+y^2+z=4$, ou seja,
  $x^2+y^2=4-z$. Trata-se de um \emph{parabolóide de revolução} com
  vértice em $(0,0,4)$, aberto para baixo.

  \textbf{Passo 2 -- Curvas de nível.} Para $z=c$:
  \[
    x^2+y^2 = 4-c.
  \]
  São circunferências de raio $\sqrt{4-c}$, definidas para $c\leq 4$.
  Quando $c=4$ obtém-se o ponto $(0,0)$; quando $c\to-\infty$ o raio
  cresce indefinidamente.

  \textbf{Passo 3 -- Cortes verticais.} Fixando $y=0$:
  $z=4-x^2$, que é uma parábola no plano $Oxz$. Por simetria de
  revolução, qualquer corte vertical pelo eixo $Oz$ dá a mesma parábola.
\end{solbox}

\begin{exbox}{3 --- Determinação de Domínio}
  Determine o domínio de $f(x,y)=\dfrac{\ln(1-x^2-y^2)}{\sqrt{x+y}}$.
\end{exbox}
\begin{solbox}
  É necessário que:
  \begin{itemize}
    \item $1-x^2-y^2>0 \Rightarrow x^2+y^2 < 1$ (disco unitário aberto).
    \item $x+y>0$ (acima da reta $y=-x$).
  \end{itemize}
  Assim $D = \{(x,y)\mid x^2+y^2<1 \text{ e } x+y>0\}$.
\end{solbox}

% ────────────────────────────────────────────────────────
\section{Limites e Continuidade}

% ────────────────────────────────────────────────────────
\subsection{Enquadramento Teórico}

\begin{defbox}[title={Limite em $\mathbb{R}^2$}]
  Diz-se que $\displaystyle\lim_{(x,y)\to(a,b)}f(x,y) = L$ se para todo o
  $\varepsilon>0$ existe $\delta>0$ tal que
  \[
    0 < \sqrt{(x-a)^2+(y-b)^2} < \delta \implies |f(x,y)-L|<\varepsilon.
  \]
\end{defbox}

\begin{notebox}[title={Estratégia: Provar que um Limite Não Existe}]
  Ao contrário de $\mathbb{R}$, em $\mathbb{R}^2$ o limite tem de ser o mesmo
  ao longo de \emph{todos} os caminhos. Uma estratégia comum para mostrar que
  um limite \emph{não existe} consiste em exibir dois caminhos que dão valores
  diferentes.
\end{notebox}

\begin{defbox}[title={Continuidade em $\mathbb{R}^2$}]
  $f$ é \emph{contínua} em $(a,b)$ se
  $\displaystyle\lim_{(x,y)\to(a,b)}f(x,y)=f(a,b)$.
\end{defbox}

% ────────────────────────────────────────────────────────
\subsection{Exemplos Resolvidos}

\begin{exbox}{4 --- Limite Inexistente}
  Mostre que $\displaystyle\lim_{(x,y)\to(0,0)}\frac{xy}{x^2+y^2}$ não existe.
\end{exbox}
\begin{solbox}
  Pelo caminho $y=0$: $f(x,0)=0\to 0$.\\
  Pelo caminho $y=x$: $f(x,x)=\dfrac{x^2}{2x^2}=\dfrac{1}{2}\to\dfrac{1}{2}$.\\
  Como os limites são diferentes, o limite não existe.
\end{solbox}

% ────────────────────────────────────────────────────────
\section{Derivadas Parciais e Diferenciabilidade}

% ────────────────────────────────────────────────────────
\subsection{Enquadramento Teórico}

\begin{defbox}[title={Derivada Parcial}]
  A \emph{derivada parcial} de $f$ em ordem a $x$ no ponto $(a,b)$ é
  \[
    \frac{\partial f}{\partial x}(a,b)
    = \lim_{h\to 0}\frac{f(a+h,b)-f(a,b)}{h},
  \]
  calculada derivando $f$ em ordem a $x$ mantendo $y$ constante.
\end{defbox}

\begin{defbox}[title={Gradiente}]
  O \emph{gradiente} de $f$ em $(a,b)$ é o vetor
  \[
    \nabla f(a,b) = \left(\frac{\partial f}{\partial x}(a,b),\,
                          \frac{\partial f}{\partial y}(a,b)\right).
  \]
  Aponta na direção de maior crescimento de $f$.
\end{defbox}

\begin{defbox}[title={Derivada Direcional}]
  Para um vetor unitário $\mathbf{u}=(u_1,u_2)$, a derivada direcional é
  \[
    D_{\mathbf{u}}f(a,b) = \nabla f(a,b)\cdot\mathbf{u}
      = \frac{\partial f}{\partial x}u_1 + \frac{\partial f}{\partial y}u_2.
  \]
\end{defbox}

\begin{defbox}[title={Diferenciabilidade}]
  $f$ é \emph{diferenciável} em $(a,b)$ se existir uma aplicação linear
  $L:\mathbb{R}^2\to\mathbb{R}$ tal que
  \[
    \lim_{(h,k)\to(0,0)}\frac{f(a+h,b+k)-f(a,b)-L(h,k)}{\sqrt{h^2+k^2}}=0.
  \]
  A aplicação linear é $L(h,k)=\nabla f(a,b)\cdot(h,k)$.
\end{defbox}

\begin{thmbox}[title={Condição Suficiente de Diferenciabilidade e Cadeia de Implicações}]
  Se $f_x$ e $f_y$ existirem e forem contínuas numa vizinhança de $(a,b)$,
  então $f$ é diferenciável em $(a,b)$.

  \medskip
  \textbf{Cadeia de implicações:}
  Diferenciável $\Rightarrow$ Contínua $\Rightarrow$ Limite existe;
  mas Derivadas parciais existem $\not\Rightarrow$ Diferenciável.
\end{thmbox}

% ────────────────────────────────────────────────────────
\subsection{Exemplos Resolvidos}

\begin{exbox}{5 --- Cálculo do Gradiente}
  Determine $\nabla f$ para $f(x,y)=x^3 y + \sin(xy)$.
\end{exbox}
\begin{solbox}
  \[
    f_x = 3x^2 y + y\cos(xy), \qquad f_y = x^3 + x\cos(xy).
  \]
\end{solbox}

\begin{exbox}{6 --- Derivadas Parciais de Segunda Ordem e Hessiana}
  Para $f(x,y)=e^{x^2+y^2}$, calcule todas as derivadas parciais de
  segunda ordem.
\end{exbox}
\begin{solbox}
  $f_x = 2xe^{x^2+y^2}$, $f_y = 2ye^{x^2+y^2}$.
  \[
    f_{xx}=(2+4x^2)e^{x^2+y^2},\quad f_{yy}=(2+4y^2)e^{x^2+y^2},\quad
    f_{xy}=f_{yx}=4xye^{x^2+y^2}.
  \]
  (Note-se que $f_{xy}=f_{yx}$ pelo Teorema de Schwarz/Clairaut, uma vez
  que todas as derivadas parciais de segunda ordem são contínuas.)
\end{solbox}


% ────────────────────────────────────────────────────────
\section{Regra da Cadeia}

% ────────────────────────────────────────────────────────
\subsection{Enquadramento Teórico}

\begin{thmbox}[title={Regra da Cadeia para Funções de Várias Variáveis}]
  Sejam $z=f(x,y)$, $x=x(t)$, $y=y(t)$ diferenciáveis. Então
  \[
    \frac{dz}{dt} = \frac{\partial f}{\partial x}\frac{dx}{dt}
                  + \frac{\partial f}{\partial y}\frac{dy}{dt}.
  \]
  De modo mais geral, se $x=x(s,t)$, $y=y(s,t)$:
  \[
    \frac{\partial z}{\partial s}
      = \frac{\partial f}{\partial x}\frac{\partial x}{\partial s}
      + \frac{\partial f}{\partial y}\frac{\partial y}{\partial s}.
  \]
\end{thmbox}

% ────────────────────────────────────────────────────────
\subsection{Exemplos Resolvidos}

\begin{exbox}{7 --- Aplicação da Regra da Cadeia}
  Seja $z = x^2+y^2$, $x=\cos t$, $y=\sin t$. Calcule $\dfrac{dz}{dt}$.
\end{exbox}
\begin{solbox}
  \[
    \frac{dz}{dt} = 2x(-\sin t)+2y(\cos t) = -2\cos t\sin t + 2\sin t\cos t = 0.
  \]
  (Faz sentido: $z = \cos^2 t + \sin^2 t = 1$, constante.)
\end{solbox}

% ────────────────────────────────────────────────────────
\section{Extremos Livres}

% ────────────────────────────────────────────────────────
\subsection{Enquadramento Teórico}

\begin{defbox}[title={Ponto Crítico}]
  $(a,b)$ é um \emph{ponto crítico} de $f$ se $\nabla f(a,b)=\mathbf{0}$,
  ou seja, $f_x(a,b)=0$ e $f_y(a,b)=0$.
\end{defbox}

\begin{thmbox}[title={Teste da Segunda Derivada (Critério de Classificação)}]
  Seja $(a,b)$ um ponto crítico de $f$. Define-se o \emph{determinante Hessiano}
  \[
    H = f_{xx}(a,b)\,f_{yy}(a,b) - [f_{xy}(a,b)]^2.
  \]
  \begin{itemize}
    \item $H > 0$ e $f_{xx}(a,b)>0$: \textbf{mínimo local}.
    \item $H > 0$ e $f_{xx}(a,b)<0$: \textbf{máximo local}.
    \item $H < 0$: \textbf{ponto de sela}.
    \item $H = 0$: o teste é \textbf{inconclusivo}.
  \end{itemize}
\end{thmbox}

% ────────────────────────────────────────────────────────
\subsection{Exemplos Resolvidos}

\begin{exbox}{8 --- Classificação de Pontos Críticos}
  Determine e classifique os pontos críticos de $f(x,y)=x^3-3x+y^2-2y$.
\end{exbox}
\begin{solbox}
  \[
    f_x = 3x^2-3 = 0 \implies x=\pm 1,\qquad
    f_y = 2y-2 = 0 \implies y = 1.
  \]
  Pontos críticos: $(1,1)$ e $(-1,1)$.

  Derivadas de segunda ordem: $f_{xx}=6x$, $f_{yy}=2$, $f_{xy}=0$.

  \begin{center}
  \renewcommand{\arraystretch}{1.3}
  \begin{tabular}{cccccc}
    \toprule
    Ponto & $f_{xx}$ & $f_{yy}$ & $f_{xy}$ & $H$ & Classificação \\
    \midrule
    $(1,1)$  & $6$  & $2$ & $0$ & $12>0$, $f_{xx}>0$ & Mínimo local \\
    $(-1,1)$ & $-6$ & $2$ & $0$ & $-12<0$ & Ponto de sela \\
    \bottomrule
  \end{tabular}
  \end{center}

  Valor mínimo local: $f(1,1) = 1-3+1-2 = -3$.
\end{solbox}

% ────────────────────────────────────────────────────────
\subsection{Exercícios Práticos}

\begin{exercisebox}{1}
  Para $f(x,y)=\sqrt{9-x^2-y^2}$, determine o domínio e descreva as
  curvas de nível.
\end{exercisebox}
\begin{solbox}
  \textbf{Passo 1 -- Domínio.} A raiz quadrada exige que o radicando seja
  não-negativo:
  \[
    9-x^2-y^2 \geq 0 \iff x^2+y^2 \leq 9.
  \]
  O domínio é o disco fechado $D = \{(x,y)\in\mathbb{R}^2 \mid x^2+y^2\leq 9\}$,
  de raio $3$ centrado na origem.

  \textbf{Passo 2 -- Curvas de nível.} Para $c\geq 0$, a equação $f(x,y)=c$
  equivale a $\sqrt{9-x^2-y^2}=c$, ou seja,
  \[
    x^2+y^2 = 9-c^2, \quad \text{com } 0\leq c\leq 3.
  \]
  \begin{itemize}
    \item Se $0\leq c < 3$: circunferência de raio $\sqrt{9-c^2}$.
    \item Se $c=3$: o ponto $(0,0)$ (raio nulo).
    \item Se $c>3$ ou $c<0$: curva vazia (fora do contradomínio).
  \end{itemize}
  As curvas de nível são portanto circunferências concêntricas de raio
  decrescente à medida que $c$ aumenta de $0$ a $3$, correspondendo
  geometricamente a paralelos de uma semiesfera de raio $3$.
\end{solbox}

\begin{exercisebox}{2}
  Averigue se os seguintes limites existem e calcule-os caso existam:
  \begin{enumerate}[label=(\alph*)]
    \item $\displaystyle\lim_{(x,y)\to(0,0)}\frac{x^2 y}{x^4+y^2}$
    \item $\displaystyle\lim_{(x,y)\to(1,2)}(x^2+3xy)$
    \item $\displaystyle\lim_{(x,y)\to(0,0)}\frac{x^3+y^3}{x^2+y^2}$
  \end{enumerate}
\end{exercisebox}
\begin{solbox}
  \textbf{(a)} Testa-se com dois caminhos distintos.

  \emph{Caminho 1:} $y=0$. $f(x,0)=\dfrac{0}{x^4}=0\to 0$.

  \emph{Caminho 2:} $y=x^2$. $f(x,x^2)=\dfrac{x^2\cdot x^2}{x^4+x^4}=\dfrac{x^4}{2x^4}=\dfrac{1}{2}\to\dfrac{1}{2}$.

  Como os dois limites são diferentes ($0\neq\tfrac{1}{2}$), o limite
  \textbf{não existe}.

  \textbf{(b)} A função $g(x,y)=x^2+3xy$ é polinomial, portanto contínua em
  todo o $\mathbb{R}^2$. Assim:
  \[
    \lim_{(x,y)\to(1,2)}(x^2+3xy) = 1^2+3\cdot 1\cdot 2 = 1+6 = 7.
  \]

  \textbf{(c)} Passa-se para coordenadas polares: $x=r\cos\theta$,
  $y=r\sin\theta$, com $r\to 0^+$.
  \[
    \frac{x^3+y^3}{x^2+y^2}
    = \frac{r^3(\cos^3\theta+\sin^3\theta)}{r^2}
    = r(\cos^3\theta+\sin^3\theta).
  \]
  Como $|\cos^3\theta+\sin^3\theta|\leq 2$ para todo o $\theta$, temos
  $\left|r(\cos^3\theta+\sin^3\theta)\right|\leq 2r\to 0$.
  Portanto o limite \textbf{existe e é igual a $0$}.
\end{solbox}

\begin{exercisebox}{3}
  Calcule todas as derivadas parciais de primeira e segunda ordem de
  $f(x,y) = x^2\sin y + y\ln x$.
\end{exercisebox}
\begin{solbox}
  \textbf{Passo 1 -- Derivadas de primeira ordem.}
  \begin{align*}
    f_x &= \frac{\partial}{\partial x}(x^2\sin y) + \frac{\partial}{\partial x}(y\ln x)
         = 2x\sin y + \frac{y}{x}.\\[4pt]
    f_y &= \frac{\partial}{\partial y}(x^2\sin y) + \frac{\partial}{\partial y}(y\ln x)
         = x^2\cos y + \ln x.
  \end{align*}

  \textbf{Passo 2 -- Derivadas de segunda ordem.}
  \begin{align*}
    f_{xx} &= \frac{\partial}{\partial x}\!\left(2x\sin y+\frac{y}{x}\right)
            = 2\sin y - \frac{y}{x^2}.\\[4pt]
    f_{yy} &= \frac{\partial}{\partial y}(x^2\cos y+\ln x)
            = -x^2\sin y.\\[4pt]
    f_{xy} &= \frac{\partial}{\partial y}\!\left(2x\sin y+\frac{y}{x}\right)
            = 2x\cos y + \frac{1}{x}.\\[4pt]
    f_{yx} &= \frac{\partial}{\partial x}(x^2\cos y+\ln x)
            = 2x\cos y + \frac{1}{x}.
  \end{align*}
  Confirma-se que $f_{xy}=f_{yx}$ (Teorema de Schwarz), pois ambas as
  derivadas mistas são contínuas no domínio $x>0$.
\end{solbox}

\begin{exercisebox}{4}
  Seja $f(x,y) = x^2 y - y^3$. Determine a derivada direcional em $(1,1)$
  na direção de $\mathbf{v}=(3,4)$.
\end{exercisebox}
\begin{solbox}
  \textbf{Passo 1 -- Normalizar o vetor.}
  \[
    \|\mathbf{v}\| = \sqrt{3^2+4^2} = 5,\qquad
    \mathbf{u} = \frac{\mathbf{v}}{\|\mathbf{v}\|} = \left(\frac{3}{5},\frac{4}{5}\right).
  \]

  \textbf{Passo 2 -- Calcular o gradiente.}
  \[
    f_x = 2xy,\quad f_y = x^2 - 3y^2.
  \]
  Em $(1,1)$: $f_x(1,1)=2$, $f_y(1,1)=1-3=-2$.
  Logo $\nabla f(1,1)=(2,-2)$.

  \textbf{Passo 3 -- Produto interno.}
  \[
    D_{\mathbf{u}}f(1,1) = \nabla f(1,1)\cdot\mathbf{u}
    = 2\cdot\frac{3}{5}+(-2)\cdot\frac{4}{5}
    = \frac{6}{5}-\frac{8}{5} = -\frac{2}{5}.
  \]
  A derivada direcional é $-\dfrac{2}{5}$, ou seja, $f$ decresce nessa
  direção.
\end{solbox}

\begin{exercisebox}{5}
  Dado $z = e^{u^2+v^2}$, $u = x\cos y$, $v = x\sin y$, utilize a regra
  da cadeia para calcular $\partial z/\partial x$ e $\partial z/\partial y$.
\end{exercisebox}
\begin{solbox}
  \textbf{Passo 1 -- Derivadas intermédias.}
  \[
    \frac{\partial z}{\partial u}=2ue^{u^2+v^2},\quad
    \frac{\partial z}{\partial v}=2ve^{u^2+v^2}.
  \]
  \[
    \frac{\partial u}{\partial x}=\cos y,\quad
    \frac{\partial u}{\partial y}=-x\sin y,\quad
    \frac{\partial v}{\partial x}=\sin y,\quad
    \frac{\partial v}{\partial y}=x\cos y.
  \]

  \textbf{Passo 2 -- Regra da cadeia para $\partial z/\partial x$.}
  \[
    \frac{\partial z}{\partial x}
    = 2ue^{u^2+v^2}\cos y + 2ve^{u^2+v^2}\sin y
    = 2e^{u^2+v^2}(u\cos y + v\sin y).
  \]

  \textbf{Passo 3 -- Simplificar.} Substituindo $u=x\cos y$ e $v=x\sin y$:
  $u\cos y+v\sin y = x\cos^2 y + x\sin^2 y = x$.
  Como $u^2+v^2=x^2$:
  \[
    \frac{\partial z}{\partial x} = 2xe^{x^2}.
  \]

  \textbf{Passo 4 -- Regra da cadeia para $\partial z/\partial y$.}
  \[
    \frac{\partial z}{\partial y}
    = 2ue^{u^2+v^2}(-x\sin y) + 2ve^{u^2+v^2}(x\cos y)
    = 2xe^{x^2}(-u\sin y+v\cos y).
  \]
  Substituindo: $-u\sin y+v\cos y = -x\cos y\sin y+x\sin y\cos y = 0$.
  \[
    \frac{\partial z}{\partial y} = 0.
  \]
  (Faz sentido: $z=e^{x^2}$ não depende de $y$.)
\end{solbox}

\begin{exercisebox}{6}
  Determine e classifique todos os pontos críticos de:
  \begin{enumerate}[label=(\alph*)]
    \item $f(x,y) = x^2 + xy + y^2 - 3x$
    \item $f(x,y) = x^3 + y^3 - 3xy$
    \item $f(x,y) = \sin x + \sin y + \sin(x+y)$, em
          $(0,\pi/2)\times(0,\pi/2)$
  \end{enumerate}
\end{exercisebox}
\begin{solbox}
  \textbf{(a)} $f(x,y) = x^2+xy+y^2-3x$.

  \emph{Condições de primeira ordem.}
  \[
    f_x = 2x+y-3=0,\qquad f_y = x+2y=0.
  \]
  De $f_y=0$: $x=-2y$. Substituindo em $f_x=0$: $-4y+y-3=0\Rightarrow y=-1$,
  $x=2$. \quad Ponto crítico: $(2,-1)$.

  \emph{Teste da segunda derivada.}
  $f_{xx}=2$, $f_{yy}=2$, $f_{xy}=1$.
  $H = 2\cdot 2 - 1^2 = 3 > 0$, $f_{xx}=2>0$.
  $\Rightarrow$ $(2,-1)$ é \textbf{mínimo local}, com $f(2,-1)=-3$.

  \bigskip
  \textbf{(b)} $f(x,y)=x^3+y^3-3xy$.

  \emph{Condições de primeira ordem.}
  $f_x = 3x^2-3y=0 \Rightarrow y=x^2$; $f_y = 3y^2-3x=0 \Rightarrow x=y^2$.
  Substituindo: $x=x^4\Rightarrow x\in\{0,1\}$,
  logo pontos críticos $(0,0)$ e $(1,1)$.

  \emph{Teste da segunda derivada.}
  $f_{xx}=6x$, $f_{yy}=6y$, $f_{xy}=-3$.

  \begin{center}
  \renewcommand{\arraystretch}{1.3}
  \begin{tabular}{cccccc}
    \toprule
    Ponto & $f_{xx}$ & $f_{yy}$ & $f_{xy}$ & $H$ & Classificação \\
    \midrule
    $(0,0)$ & $0$ & $0$ & $-3$ & $-9<0$ & Ponto de sela \\
    $(1,1)$ & $6$ & $6$ & $-3$ & $27>0$, $f_{xx}>0$ & Mínimo local \\
    \bottomrule
  \end{tabular}
  \end{center}

  Valor mínimo local: $f(1,1)=1+1-3=-1$.

  \bigskip
  \textbf{(c)} $f(x,y)=\sin x+\sin y+\sin(x+y)$, em $(0,\tfrac{\pi}{2})\times(0,\tfrac{\pi}{2})$.

  \emph{Condições de primeira ordem.}
  \[
    f_x = \cos x + \cos(x+y) = 0,\qquad
    f_y = \cos y + \cos(x+y) = 0.
  \]
  Subtraindo: $\cos x = \cos y$. No domínio $(0,\tfrac{\pi}{2})$ o cosseno é estritamente
  decrescente, logo $x=y$.

  Substituindo $y=x$ em $f_x=0$:
  $\cos x + \cos 2x = 0$.
  Usando $\cos 2x = 2\cos^2 x - 1$:
  $(2\cos x - 1)(\cos x + 1) = 0$.
  No domínio: $\cos x = \tfrac{1}{2} \Rightarrow x=\tfrac{\pi}{3}$.
  Ponto crítico: $\left(\tfrac{\pi}{3},\tfrac{\pi}{3}\right)$.

  \emph{Teste da segunda derivada.}
  Em $\left(\tfrac{\pi}{3},\tfrac{\pi}{3}\right)$: $\sin\tfrac{\pi}{3}=\sin\tfrac{2\pi}{3}=\tfrac{\sqrt{3}}{2}$.
  \[
    f_{xx}=f_{yy}=-\sqrt{3},\quad f_{xy}=-\frac{\sqrt{3}}{2},\quad
    H = 3-\frac{3}{4}=\frac{9}{4}>0,\; f_{xx}<0.
  \]
  $\Rightarrow$ $\left(\tfrac{\pi}{3},\tfrac{\pi}{3}\right)$ é \textbf{máximo local},
  com $f=\tfrac{3\sqrt{3}}{2}$.
\end{solbox}

% ────────────────────────────────────────────────────────
\subsection{Questões Propostas para Defesa Oral}

\begin{oralbox}
  \textbf{Q1.} Explique por que razão a existência de derivadas parciais num ponto
  \emph{não} garante a continuidade nesse ponto. Dê um contra-exemplo.
\end{oralbox}

\begin{oralbox}
  \textbf{Q2.} Enuncie as condições em que as derivadas parciais mistas $f_{xy}$
  e $f_{yx}$ são iguais (Teorema de Schwarz/Clairaut).
\end{oralbox}

\begin{oralbox}
  \textbf{Q3.} Interprete geometricamente o gradiente: o que nos diz $\nabla f(a,b)$
  sobre a superfície $z=f(x,y)$?
\end{oralbox}

\begin{oralbox}
  \textbf{Q4.} Deduza a fórmula da regra da cadeia para $dz/dt$ quando
  $z=f(x(t),y(t))$ a partir da definição de diferenciabilidade.
\end{oralbox}

\begin{oralbox}
  \textbf{Q5.} Descreva o teste da segunda derivada. O que acontece quando o
  determinante Hessiano é zero? Apresente exemplos de cada caso.
\end{oralbox}

\begin{oralbox}
  \textbf{Q6.} Como generaliza o conceito de derivada direcional a noção de
  derivada parcial? Quando é que $D_{\mathbf{u}}f$ é maximizada, e em
  que direção?
\end{oralbox}

\begin{oralbox}
  \textbf{Q7.} Explique a ligação entre curvas de nível, o gradiente e a direção
  de maior subida/descida.
\end{oralbox}

% ---------------------------------------------------------
% CAPÍTULO 4
% ---------------------------------------------------------
% ════════════════════════════════════════════════════════════════
% CAPÍTULO 4 -- INTEGRAIS DUPLOS E MODELAÇÃO COMPUTACIONAL
% Ficheiro para \include{} no template principal.
%
% NOTA PARA O TEMPLATE PRINCIPAL: Este capítulo utiliza código
% MATLAB/Octave. Adicionar ao preâmbulo do ficheiro principal:
%
%   \usepackage{listings}
%   \lstset{
%       language=Matlab,
%       basicstyle=\ttfamily\footnotesize,
%       keywordstyle=\color{blue}\bfseries,
%       commentstyle=\color{teal},
%       stringstyle=\color{red},
%       numbers=left,
%       numberstyle=\tiny\color{gray},
%       stepnumber=1,
%       breaklines=true,
%       frame=single,
%       backgroundcolor=\color{gray!5},
%       captionpos=b,
%       showstringspaces=false
%   }
% ════════════════════════════════════════════════════════════════

\chapter{Integrais Duplos e Modelação Computacional}

% ────────────────────────────────────────────────────────
\section{Fundamentos Teóricos}

% ────────────────────────────────────────────────────────
\subsection{O Conceito de Integral Duplo}

Enquanto um integral simples calcula a área debaixo de uma curva num gráfico 2D, um integral duplo calcula o volume abaixo de uma superfície num espaço 3D. A notação padrão é dada por:

\begin{defbox}[title={Integral Duplo}]
\[
    V = \iint_R f(x,y) \,\dd A
\]
onde $R$ representa a região planar base (o domínio) e $f(x,y)$ a função que dita a altura da superfície em cada ponto $(x,y)$.
\end{defbox}

% ────────────────────────────────────────────────────────
\subsection{Definição Formal e Somas de Riemann}

A construção rigorosa de um integral duplo baseia-se na expansão das Somas de Riemann. Seja $f(x,y)$ uma função definida numa região planar fechada e limitada $R$. Se dividirmos $R$ numa malha de pequenos retângulos de área $\Delta A$, e escolhermos um ponto de amostragem $(x_i^*, y_j^*)$ dentro de cada sub-retângulo, o volume aproximado do sólido abaixo da superfície é a soma dos volumes de vários paralelepípedos:

\begin{defbox}[title={Soma de Riemann Dupla}]
\[
    V \approx \sum_{i=1}^{n} \sum_{j=1}^{m} f(x_i^*, y_j^*) \Delta A
\]
O integral duplo é definido como o limite desta soma quando a área dos sub-retângulos tende para zero (ou seja, quando a malha se torna infinitamente fina):
\[
    \iint_R f(x,y) \,\dd A = \lim_{\Delta A \to 0} \sum_{i=1}^{n} \sum_{j=1}^{m} f(x_i^*, y_j^*) \Delta A
\]
\end{defbox}

\begin{notebox}
Computacionalmente, os métodos de integração numérica (como a regra dos trapézios bidimensional usada nos projetos) são aproximações algorítmicas diretas desta mesma Soma de Riemann.
\end{notebox}

% ────────────────────────────────────────────────────────
\subsection{Teorema de Fubini, Integração Iterada e Tipos de Regiões}

O cálculo direto pelo limite da soma de Riemann é impraticável. O Teorema de Fubini é a ponte que permite calcular integrais duplos transformando-os em dois integrais simples sucessivos. Este teorema resolve o problema afirmando que, se a função $f(x,y)$ for contínua na região $R$, o integral duplo pode ser avaliado através de integrais iterados unidimensionais. A complexidade do cálculo depende da geometria da região $R$:

\begin{thmbox}[title={Teorema de Fubini e Integração Iterada}]
\begin{itemize}
    \item \textbf{Regiões do Tipo I (Simples na vertical):} A região é limitada à esquerda e à direita por constantes $x=a$ e $x=b$, e em cima e em baixo por funções contínuas $y=g_1(x)$ e $y=g_2(x)$. Integra-se primeiro em ordem a $y$:
    \[
        \iint_R f(x,y) \,\dd A = \int_{a}^{b} \left[ \int_{g_1(x)}^{g_2(x)} f(x,y) \,\dd y \right] \dd x
    \]

    \item \textbf{Regiões do Tipo II (Simples na horizontal):} A região é limitada em cima e em baixo por constantes $y=c$ e $y=d$, e lateralmente por funções $x=h_1(y)$ e $x=h_2(y)$. Integra-se primeiro em ordem a $x$:
    \[
        \iint_R f(x,y) \,\dd A = \int_{c}^{d} \left[ \int_{h_1(y)}^{h_2(y)} f(x,y) \,\dd x \right] \dd y
    \]
\end{itemize}
\end{thmbox}

\begin{tipbox}
Dependendo da geometria da região $R$, a capacidade de alterar a ordem de integração (de Tipo I --- $\dd x\,\dd y$ para Tipo II --- $\dd y\,\dd x$) é uma competência analítica vital, pois muitas vezes um integral iterado é impossível de resolver analiticamente numa ordem, mas trivial na outra, podendo este processo simplificar drasticamente a resolução do problema.
\end{tipbox}

% ────────────────────────────────────────────────────────
\subsection{Transformação para Coordenadas Polares}

Quando a região de integração possui simetria circular ou a função integranda contém termos como $x^2 + y^2$, o uso de coordenadas cartesianas $(x,y)$ gera fronteiras de integração complexas com raízes quadradas. Nestes casos, aplica-se uma transformação geométrica para o sistema de coordenadas polares $(r, \theta)$, onde:
\[
    x = r \cos\theta \quad \text{e} \quad y = r \sin\theta
\]

\begin{thmbox}[title={Mudança para Coordenadas Polares e Jacobiano}]
Ao deformarmos o espaço bidimensional para um novo sistema de coordenadas, a área dos elementos infinitesimais deforma-se. O fator de escala dessa deformação é dado pelo valor absoluto do determinante Jacobiano da transformação. Para coordenadas polares, o Jacobiano $J$ calcula-se como:
\[
    J = \begin{vmatrix}
    \dfrac{\partial x}{\partial r} & \dfrac{\partial x}{\partial \theta} \\[6pt]
    \dfrac{\partial y}{\partial r} & \dfrac{\partial y}{\partial \theta}
    \end{vmatrix}
    = \begin{vmatrix}
    \cos\theta & -r\sin\theta \\
    \sin\theta & r\cos\theta
    \end{vmatrix}
    = r(\cos^2\theta + \sin^2\theta) = r
\]
Portanto, o elemento de área diferencial muda de $\dd x\,\dd y$ para $r \,\dd r \,\dd\theta$. O teorema geral da mudança de variáveis estabelece que:
\[
    \iint_R f(x,y) \,\dd x\,\dd y = \iint_{S} f(r\cos\theta, r\sin\theta) \, r \,\dd r\,\dd\theta
\]
onde $S$ é a representação da região $R$ no novo plano $r\theta$.

Em suma:
\[
    \dd A = \dd x \, \dd y = r \, \dd r \, \dd\theta
\]
\end{thmbox}

% ────────────────────────────────────────────────────────
\section{Modelação Computacional e Projetos Interdisciplinares}

A verdadeira avaliação das competências ocorre quando os conceitos analíticos são traduzidos em algoritmos. Computadores resolvem integrais através de métodos numéricos (somatórios em malhas discretas). Abaixo detalham-se os três projetos de aplicação prática.

% ────────────────────────────────────────────────────────
\subsection{Projeto 1: Simulação de Difusão de Calor}

A energia térmica numa superfície planar acumula-se e difunde-se ao longo do tempo. Para calcular o calor total acumulado, integra-se a função de temperatura sobre a área 2D.

\begin{lstlisting}[caption={Simulação de Difusão de Calor em MATLAB/Octave}]
% 1. Configuracao do Dominio (Regiao Planar R)
nx = 50; ny = 50; dx = 0.1; dy = 0.1;
T = zeros(nx, ny);
T(20:30, 20:30) = 100; % Fonte de calor inicial

% 2. Simulacao da Difusao (Diferencas Finitas)
alpha = 0.2;
for k = 1:50
    T_nova = T;
    for i = 2:nx-1
        for j = 2:ny-1
            T_nova(i,j) = T(i,j) + alpha * (T(i+1,j) + T(i-1,j) + T(i,j+1) + T(i,j-1) - 4*T(i,j));
        end
    end
    T = T_nova;
end

% 3. Calculo do Integral Duplo Numerico (Energia Total)
calor_total = sum(sum(T)) * (dx * dy);

% 4. Visualizacao 3D da Superficie
surf(T); colormap('hot');
title(sprintf('Energia Total: %.2f', calor_total));
\end{lstlisting}

% ────────────────────────────────────────────────────────
\subsection{Projeto 2: Análise do Movimento Solar e Vetores 3D}

Para estimar a radiação total capturada por um painel, modela-se o painel como uma região $R$ no espaço e calcula-se o ângulo de incidência do vetor solar. A energia total diária exige integração dupla sobre a área do painel iterada ao longo do tempo.

\begin{lstlisting}[caption={Análise de Incidência Solar em MATLAB/Octave}]
x = 0:0.1:2; y = 0:0.1:1.5; % Dimensoes do painel
[X, Y] = meshgrid(x, y);
normal_painel = [0, 0, 1];
energia_total = 0;

for t = 0:1:12
    theta = (pi/12) * t;
    vetor_sol = [cos(theta), 0, sin(theta)];

    % Intensidade via produto interno (Vetores 3D)
    intensidade = max(0, dot(normal_painel, vetor_sol));

    % Funcao densidade de energia f(x,y) (ex: lado sombreado)
    Sombra = 1 - 0.1 * X;
    Energia_XY = intensidade * Sombra;

    % Integral Duplo iterado (Metodo dos Trapezios)
    energia_instante = trapz(y, trapz(x, Energia_XY, 2));
    energia_total = energia_total + energia_instante;
end
\end{lstlisting}

% ────────────────────────────────────────────────────────
\subsection{Projeto 3: Redes Neurais e Descida de Gradiente}

Ao treinar uma rede neural, otimiza-se uma função de custo $L(w_1, w_2)$ que representa uma superfície no espaço dos parâmetros. A compreensão das derivadas parciais e da topologia de funções de várias variáveis é crucial para o algoritmo de \textit{Gradient Descent}.

\begin{lstlisting}[caption={Otimização em Superfície de Erro (Redes Neurais)}]
% 1. Espaco de Parametros (Pesos w1, w2)
[w1, w2] = meshgrid(-5:0.2:5, -5:0.2:5);
Custo = w1.^2 + 2 * w2.^2; % Funcao de Custo (Superficie 3D)

% 2. Otimizacao via Descida de Gradiente
peso = [4, 4]; % Ponto inicial
taxa = 0.1;
hist_pesos = zeros(21, 2); hist_pesos(1,:) = peso;

for i = 1:20
    grad = [2*peso(1), 4*peso(2)]; % Derivadas parciais
    peso = peso - taxa * grad;
    hist_pesos(i+1,:) = peso;
end

% 3. Curvas de Nivel no Plano 2D
contour(w1, w2, Custo, 20); hold on;
plot(hist_pesos(:,1), hist_pesos(:,2), 'r-o', 'LineWidth', 2);
title('Descida de Gradiente na Superficie de Custo');
\end{lstlisting}

% ---------------------------------------------------------
% CAPÍTULO 5
% ---------------------------------------------------------
% ════════════════════════════════════════════════════════════════
% CAPÍTULO 5 - EDOs DE PRIMEIRA ORDEM
% ════════════════════════════════════════════════════════════════

\chapter{Equações Diferenciais Ordinárias (EDO) de 1ª Ordem}

% ────────────────────────────────────────────────────────
\section{Enquadramento Teórico}

\subsection{Introdução às Equações Diferenciais}

Uma \textbf{equação diferencial} é uma equação que relaciona uma função desconhecida com uma ou mais das suas derivadas. As equações diferenciais são a linguagem natural das ciências físicas: descrevem como os sistemas variam no tempo ou no espaço.

\begin{defbox}[title={EDO de Primeira Ordem}]
  Uma \textbf{equação diferencial ordinária (EDO) de primeira ordem} é uma equação da forma
  \[
    F\!\bigl(x,\, y,\, y'\bigr) = 0,
    \qquad\text{ou equivalentemente}\qquad
    \frac{\dd y}{\dd x} = f(x, y),
  \]
  onde $y = y(x)$ é uma função desconhecida de uma única variável independente $x$, e $y' = \frac{\dd y}{\dd x}$ é a sua derivada primeira.
\end{defbox}

\begin{notebox}
  \textbf{Importância Prática:} \\
  \textbf{Em física:} Lei do arrefecimento de Newton, desintegração radioativa, carga de circuito RC, escoamento de fluidos (lei de Torricelli).\\
  \textbf{Em engenharia/computação:} Dinâmica populacional, gradiente descendente em Machine Learning, difusão de calor (1-D), cinética de reações químicas.
\end{notebox}

% Campo de direções
\begin{center}
  \begin{tikzpicture}
    \begin{axis}[
      width=10cm, height=6cm,
      xlabel={$x$}, ylabel={$y$},
      xmin=-0.2, xmax=4.5,
      ymin=-0.3, ymax=1.8,
      axis lines=center,
      tick style={draw=none},
      xticklabels={}, yticklabels={},
      title={\small\color{primary}\textbf{Campo de Direções Exemplo:} $\dfrac{\dd y}{\dd x} = y(1-y)$},
    ]
      % Linhas de equilíbrio
      \addplot[color=primary, thick, domain=0:4.4] {1} node[right, font=\tiny] {$y=1$};
      \addplot[color=accent, thick, domain=0:4.4] {0} node[right, font=\tiny] {$y=0$};
      % Curva logística por baixo
      \addplot[color=primary!60!black, very thick, domain=0:4.4, samples=80] { 1/(1 + 19*exp(-x)) };
      % Solução por cima
      \addplot[color=gray!70, thick, domain=0:3.0, samples=80] { 1/(1 - (1/3)*exp(-x)) };
    \end{axis}
  \end{tikzpicture}
\end{center}


\subsection{Classificação das EDOs de Primeira Ordem}

\begin{defbox}[title={Ordem e Grau}]
  \begin{itemize}
    \item A \textbf{ordem} de uma EDO é a ordem da derivada mais elevada que nela aparece.
    \item O \textbf{grau} é o expoente da derivada de maior ordem (após eliminar frações e radicais).
  \end{itemize}
\end{defbox}

\begin{defbox}[title={Linearidade e Autonomia}]
  \begin{itemize}
    \item \textbf{Linear:} Uma EDO é linear se puder ser escrita na forma $y' + P(x)\,y = Q(x)$. É \textbf{não linear} nos restantes casos.
    \item \textbf{Autónoma:} Uma EDO é autónoma se não depender explicitamente de $x$, ou seja, $y' = f(y)$. Admitem \textbf{soluções de equilíbrio} obtidas fazendo $f(y^{*})=0$.
  \end{itemize}
\end{defbox}

\begin{center}
\renewcommand{\arraystretch}{1.3}
\begin{tabularx}{0.95\linewidth}{l X X}
  \toprule
  \textbf{Categoria} & \textbf{Forma Padrão} & \textbf{Característica Principal} \\
  \midrule
  \textbf{Separável} & $g(y)\,\dd y = h(x)\,\dd x$ & Variáveis desacoplam completamente \\
  \textbf{Linear}    & $y' + P(x)y = Q(x)$ & Resolvida por fator integrante \\
  \textbf{Exata}     & $M\,\dd x + N\,\dd y = 0$, $M_y=N_x$ & Existe função potencial \\
  \textbf{Bernoulli} & $y' + P(x)y = Q(x)y^n$ & Linearizada pela substituição $v=y^{1-n}$ \\
  \textbf{Autónoma}  & $y' = f(y)$ & Sem dependência explícita em $x$ \\
  \bottomrule
\end{tabularx}
\end{center}

\subsection{Soluções: Geral e Particular}

\begin{defbox}[title={Tipos de Soluções}]
  \begin{itemize}
    \item A \textbf{solução geral} contém uma constante arbitrária $C$.
    \item A \textbf{solução particular} obtém-se especificando uma \textbf{condição inicial} $y(x_0) = y_0$.
    \item Uma \textbf{solução singular} não pode ser obtida da solução geral.
  \end{itemize}
\end{defbox}

\begin{thmbox}[title={Existência e Unicidade (Picard--Lindelöf)}]
  Se $f(x,y)$ e $\frac{\partial f}{\partial y}$ forem contínuas num retângulo contendo $(x_0, y_0)$, então o problema de valor inicial $\frac{\dd y}{\dd x} = f(x,y), \; y(x_0)=y_0$ tem uma solução \textbf{única}.
\end{thmbox}

\begin{warnbox}
  \textbf{Aviso de Existência Global:} Não assuma que a solução existe globalmente. O intervalo de existência pode ser finito. Ex: $y' = y^2, y(0)=1 \implies y = 1/(1-x)$, que diverge em $x=1$.
\end{warnbox}

\subsection{Métodos Analíticos de Resolução}

\subsubsection*{1. Equações Separáveis}
Uma EDO é separável se $\frac{\dd y}{\dd x} = g(x)h(y)$. 
\begin{enumerate}
  \item Reescrever como $\frac{\dd y}{h(y)} = g(x)\,\dd x$.
  \item Integrar: $\int \frac{\dd y}{h(y)} = \int g(x)\,\dd x + C$.
\end{enumerate}

\begin{tipbox}
  Verifique sempre $h(y)=0$ para soluções de equilíbrio antes de separar, pois dividir por $h(y)$ exige $h(y)\neq 0$.
\end{tipbox}

\subsubsection*{2. Equações Lineares (Fator Integrante)}
\begin{thmbox}[title={Fator Integrante}]
  Para $y' + P(x)y = Q(x)$, multiplicamos por $\mu(x) = e^{\int P(x)\,\dd x}$:
  \[
    \frac{\dd}{\dd x}\!\bigl[\mu(x)\,y\bigr] = \mu(x)\,Q(x) \implies y = \frac{1}{\mu(x)}\left[\int \mu(x)\,Q(x)\,\dd x + C\right].
  \]
\end{thmbox}

\subsection{Aplicações em Engenharia e Física}
\begin{itemize}
    \item \textbf{Arrefecimento de Newton:} $\frac{\dd T}{\dd t} = -k(T - T_a)$.
    \item \textbf{Circuito Elétrico RC:} $RC\,\frac{\dd V}{\dd t} + V = V_s$.
    \item \textbf{Crescimento Logístico:} $\frac{\dd P}{\dd t} = rP\left(1-\frac{P}{K}\right)$.
\end{itemize}

\subsection{Estratégias de Verificação e Erros Comuns}

\begin{tipbox}[title={Como Verificar uma Solução}]
  1. Diferenciar $y(x)$ para obter $y'(x)$. \\
  2. Substituir na EDO e verificar a igualdade. \\
  3. Verificar a condição inicial.
\end{tipbox}

\begin{warnbox}
  \textbf{Erros Frequentes:} Esquecer a constante $C$ na integração e aplicar a condição inicial antes de ter a solução geral completa.
\end{warnbox}

% ────────────────────────────────────────────────────────
\section{Exemplos Resolvidos}

\begin{exbox}{1 --- EDO Separável Básica}
  Resolver $\dfrac{\dd y}{\dd x} = xy$.
\end{exbox}
\begin{solbox}
  \textbf{Passo 1:} $y=0$ é solução. \\
  \textbf{Passo 2:} $\frac{\dd y}{y} = x\,\dd x \implies \ln|y| = \frac{x^2}{2} + C$. \\
  \textbf{Solução geral:} $\boxed{y = A\,e^{x^2/2}, \quad A\in\mathbb{R}.}$
\end{solbox}

\begin{exbox}{2 --- PVI}
  Resolver $\dfrac{\dd y}{\dd x} = -2xy^2$, com $y(0) = 1$.
\end{exbox}
\begin{solbox}
  Separar: $\frac{\dd y}{y^2} = -2x\,\dd x \implies -\frac{1}{y} = -x^2 + C$. \\
  Com $y(0)=1 \implies -1 = 0 + C \implies C = -1$. \\
  Solução particular: $\boxed{y = \frac{1}{x^2+1}.}$
\end{solbox}

\begin{exbox}{3 --- Solução Implícita}
  Resolver $\dfrac{\dd y}{\dd x} = \dfrac{x^2}{1-y^2}$.
\end{exbox}
\begin{solbox}
  $(1-y^2)\,\dd y = x^2\,\dd x \implies y - \frac{y^3}{3} = \frac{x^3}{3} + C$. \\
  Solução implícita: $\boxed{3y - y^3 = x^3 + K.}$
\end{solbox}

\begin{exbox}{4 --- Fator Integrante}
  Resolver $\dfrac{\dd y}{\dd x} - \dfrac{2}{x}\,y = x^2$.
\end{exbox}
\begin{solbox}
  $\mu = e^{\int -2/x\,\dd x} = x^{-2}$. \\
  $x^{-2}y' - 2x^{-3}y = 1 \implies \frac{\dd}{\dd x}[x^{-2}y] = 1$. \\
  Integrando: $x^{-2}y = x + C \implies \boxed{y = x^3 + Cx^2.}$
\end{solbox}

\begin{exbox}{5 --- Arrefecimento de Newton}
  Componente a $200^\circ\text{C}$ em ambiente a $25^\circ\text{C}$. Após 10 min está a $120^\circ\text{C}$.
\end{exbox}
\begin{solbox}
  $T(t) = 25 + 175e^{-kt}$. De $T(10)=120 \implies 95 = 175e^{-10k} \implies k \approx 0,061$. \\
  Para $T=50^\circ\text{C}$: $25 = 175e^{-0,061t} \implies t \approx 31,8 \text{ min}$.
\end{solbox}

\begin{exbox}{6 --- Circuito RC}
  $R=1\text{k}\Omega, C=1\text{mF}, V_s=5\text{V}, V(0)=0$.
\end{exbox}
\begin{solbox}
  EDO: $V' + V = 5$. Solução: $V(t) = 5(1-e^{-t})\text{V}$.
\end{solbox}

% ────────────────────────────────────────────────────────
\section{Exercícios Práticos}

\begin{exercisebox}{1 --- Separável}
  Resolver $\dfrac{\dd y}{\dd x} = \dfrac{y}{x}$, com $y(1)=3$.
\end{exercisebox}
\begin{solbox}
  $\frac{\dd y}{y} = \frac{\dd x}{x} \implies y = Ax$. Com $y(1)=3 \implies A=3$. Solução: $\boxed{y=3x.}$
\end{solbox}

\begin{exercisebox}{2 --- Logística}
  Resolver $\dfrac{\dd P}{\dd t} = 0,2P(1 - P/500)$, com $P(0) = 50$.
\end{exercisebox}
\begin{solbox}
  Solução da logística: $P(t) = \frac{500}{1+9e^{-0,2t}}$.
\end{solbox}

\begin{exercisebox}{3 --- Velocidade Terminal}
  $m\frac{\dd v}{\dd t} = mg - kv, \; v(0)=0$.
\end{exercisebox}
\begin{solbox}
  $v(t) = \frac{mg}{k}(1-e^{-kt/m})$.
\end{solbox}

\begin{exercisebox}{4 --- Bernoulli}
  Resolver $y' + y = y^3$.
\end{exercisebox}
\begin{solbox}
  Substituição $v = y^{-2}$. Resulta em $\boxed{y = \pm\frac{1}{\sqrt{1+Ce^{2x}}}.}$
\end{solbox}

\begin{exercisebox}{5 --- Bateria Adicional}
  a) $\frac{\dd y}{\dd x} = (1+y^2)\tan x, y(0)=1$. \\
  b) $y' + y/x = \cos x, y(\pi)=0$.
\end{exercisebox}
\begin{solbox}
  a) $y = \tan(-\ln|\cos x| + \pi/4)$. \\
  b) $y = \frac{\sin x}{x}$.
\end{solbox}

% ────────────────────────────────────────────────────────
\section{Questões para Defesa Oral}

\begin{oralbox}
  \textbf{Q1.} Enunciar o teorema de existência e unicidade. Dar um exemplo onde a unicidade falha (ex: $y' = y^{1/3}, y(0)=0$).
\end{oralbox}

\begin{oralbox}
  \textbf{Q2.} Explicar o papel do fator integrante $\mu(x)$ na transformação da EDO linear numa derivada exata.
\end{oralbox}

\begin{oralbox}
  \textbf{Q3.} Como determinar a estabilidade de equilíbrios numa EDO autónoma $y' = f(y)$ usando apenas o sinal de $f(y)$?
\end{oralbox}

\begin{oralbox}
  \textbf{Q4.} Formular a EDO de um tanque onde entra água pura e sai mistura (mistura perfeita).
\end{oralbox}

\end{document}